\documentclass[10pt,a4paper]{article}
\usepackage[margin=2.2cm]{geometry}
\usepackage{fontspec}
\setmonofont{DejaVu Sans Mono}[Scale=0.84]
\newfontfamily\unifont{DejaVu Serif}[Scale=0.9]
\usepackage{amsmath,amssymb,mathtools}
\setcounter{MaxMatrixCols}{20}
\usepackage{longtable,booktabs,array}
\usepackage{fancyvrb}
\usepackage[dvipsnames]{xcolor}
\usepackage[colorlinks=true,linkcolor=NavyBlue,urlcolor=NavyBlue]{hyperref}
\usepackage{microtype}
\usepackage{needspace}
\usepackage{tocloft}\setlength{\cftbeforesecskip}{0.45em}
\setlength{\parindent}{0pt}\setlength{\parskip}{0.55em}
\setlength{\emergencystretch}{8em}
\makeatletter
\renewcommand\section{\@startsection{section}{1}{\z@}{-3.5ex \@plus -1ex \@minus -.2ex}{2.3ex \@plus.2ex}{\normalfont\Large\bfseries\raggedright}}
\renewcommand\subsection{\@startsection{subsection}{2}{\z@}{-3.25ex\@plus -1ex \@minus -.2ex}{1.5ex \@plus .2ex}{\normalfont\large\bfseries\raggedright}}
\renewcommand\subsubsection{\@startsection{subsubsection}{3}{\z@}{-3.25ex\@plus -1ex \@minus -.2ex}{1.5ex \@plus .2ex}{\normalfont\normalsize\bfseries\raggedright}}
\makeatother
\renewcommand{\arraystretch}{1.15}
\allowdisplaybreaks
\begin{document}
\title{The world brane of 7+1 spacetime: provenance\\[0.4em]\large The task, the decisions, the inputs, the data, the derivations and the verification record of the document on the world brane}
\author{Repository: github.com/once-ere/7plus1 — prepared for Nash\\ Notebook: 7+1.nb (Wolfram Language 15.0.1)\\ Document generated by tools/make\_docs.py on 2026-10-28}
\date{}
\maketitle
\begin{abstract}
This record states what was asked, how it was interpreted, every decision taken (among them the new sections 25 and 26 of 7+1.nb, which compute the worked examples), which results of the earlier sections and which data were used and where they come from, how each result was derived, and how it was checked: the 27 checks of section 25 (27 True) and the 33 of section 26 (33 True), the 38 and 31 boolean checks of the independent Python programs tools/ref\_torus\_brane.py and tools/ref\_radion.py (38 and 31 True), the comparison of the 17 and 15 groups of numbers printed by the notebook and by them, the check that the earlier sections of the notebook are unchanged, the program versions and the SHA-256 hashes of the files.
\end{abstract}
\tableofcontents
\Needspace{12\baselineskip}\section*{1. The task}\addcontentsline{toc}{section}{1. The task}
On 28 October 2026 the author asked for the following (summarized; the original wording is the author’s private working material and is not reproduced):

\begin{enumerate}
\item Define and explain the "world brane" precisely and correctly, in its own documentation file with its own symbol file and provenance file, each in Markdown and LaTeX with the LaTeX compiled to PDF; explain it for students, clearly, with examples, well organized; do the same, separately, for the "auxiliary torus"; push the files and any other relevant files to github.com/once-ere/7plus1.
\item Then: read a note of the author on the radion (a private PDF in the workspace, not part of the repository), use its ideas to verify and clarify the radion and brane documents in preparation, and create examples in which the four extra space dimensions are (1) a 4-sphere and (2) a 4-torus. For this document that means: the brane on a 4-sphere compared with the brane on the torus (section 12 of the document), branes in a warped bulk as the note describes them (section 13), and the mixing of the Higgs boson on the brane with the radion (section 7.5).
\end{enumerate}
The standing rules of the project apply: every documentation file has a companion provenance file and a symbol file, in Markdown and LaTeX, compiled to PDF; no file refers the reader to another documentation, provenance or symbol file; every equation that is created or used is written in the file that uses it. The repository’s .gitignore is respected.

\textbf{Interpretation.} The world brane is the submanifold $W=\{x_4=x_5=x_6=x_7=0\}$ on which 7+1.nb puts the Standard Model (sections 17–24). "Precisely" is read as: a definition by an embedding, the induced metric, the extrinsic curvature, the energy–momentum tensor and the couplings, each stated as an equation; "correctly" as: every worked example is computed and checked twice, by the notebook and by an independent program. The coordinates of 7+1.nb are kept: $X=(x_0,\dots,x_7)$, $x_0$ the time (the direction of the octonion unit $e_0=1$), $x_1,x_2,x_3$ ordinary space, $x_4,\dots,x_7$ the four auxiliary space-like coordinates.

\Needspace{8\baselineskip}\subsection*{1.1 The signature: the author’s explanation of 27 October 2026}\addcontentsline{toc}{subsection}{1.1 The signature: the author’s explanation of 27 October 2026}
\textbf{Source.} On 27 October 2026 the author of the project explained that the xAct list $\{7,1,0\}$ declares seven space-like dimensions and one time-like dimension, while $\mathrm{diag}(+1,-1,\dots,-1)$ is simply the squares of the octonion basis, and asked that this be made clear in all documentation, examples and provenance records. \textbf{What it corrected.} The version of 7+1.nb pushed in the commits 52d8729 to dc2f178 of github.com/once-ere/7plus1 declared the xAct metric of section 12 with \texttt{DefMetric[\{1, 7, 0\}, …]}, reading the list as the signs of $\eta$; since then section 12 declares \texttt{DefMetric[\{7, 1, 0\}, g71[-a, -b], CD, …]}, exactly as the author’s notebook Einstein-Lovelock-7+1-Nash.nb does. \textbf{What was added.} Checks in 7+1.nb — section 4.7: $\eta$ is the squares of the octonion basis, $e_Ae_A=\eta_{AA}e_0$ and $\mathrm{Re}(e_Ae_B)=\eta_{AB}$; section 12.2: \texttt{SignatureOfMetric[g71]} $=\{7,1,0\}$ with \texttt{SignDetOfMetric[g71]} $=-1$, $\mathrm{Re}(e_Ae_B)=\eta_{AB}$, and $\mathrm{Re}(x\,x)=\eta_{AB}x^Ax^B$; section 17: the flat metric $\mathrm{diag}(+1,-1,-1,-1)$ of the world brane is the squares of the basis $e_0,e_1,e_2,e_3$ of the quaternion subalgebra — and the same statements in the Python programs \texttt{tools/\allowbreak{}ref71.py}, \texttt{tools/\allowbreak{}ref\_\allowbreak{}gravity.py} and \texttt{tools/\allowbreak{}ref\_\allowbreak{}sm.py}; the explanation itself in the text of 7+1.nb (sections 1.2, 1.4, 4.2, 4.7, 7, 11, 12.2, 13, 15.1, 17) and in every document, symbol file and provenance file of the repository. \textbf{What did not change.} No result. Compared output by output with the backup below, every output of 7+1.nb is the same, except the declared list itself, the new checks and the check counts (and, as in every run, the name of the private context and the measured times); in particular every component of every Lovelock tensor is unchanged. The components follow from the component matrix of the metric, and the sign of the determinant, \texttt{SignDetOfMetric}, is $-1$ for both lists. \textbf{Backup.} The evaluated notebook as it was before is kept in the workspace as \texttt{backups/\allowbreak{}7+1\_\allowbreak{}2026-10-27\_\allowbreak{}before\_\allowbreak{}signature.nb} (the folder backups/ is excluded from the repository by .gitignore), 1034401 bytes, SHA-256 \texttt{06094248\allowbreak{}8caf4f45\allowbreak{}120ffd37\allowbreak{}be0123e9\allowbreak{}6992a31c\allowbreak{}d7fdbde0\allowbreak{}9ee4dc4c\allowbreak{}ed4f5d3a}; \texttt{tools/\allowbreak{}verify\_\allowbreak{}7p1\_\allowbreak{}nb.py} compares every input, text and section cell of 7+1.nb with it and finds every cell present, each one unchanged or changed by this explanation, and it compares every output in the same way.

\textbf{The explanation, in full.} The 7+1 spacetime has one time-like dimension ($x_0$, the direction of the octonion unit $e_0=1$) and seven space-like dimensions ($x_1,x_2,x_3$ ordinary space, $x_4,\dots,x_7$ auxiliary space). Two different objects describe this, and neither may be replaced by the other. (i) The tensor package xAct, used in section 12 of 7+1.nb exactly as in the author’s notebook Einstein-Lovelock-7+1-Nash.nb, declares the metric with \texttt{DefMetric[\{7, 1, 0\}, g71[-a, -b], CD, …]}: the list $\{7,1,0\}$ declares seven space-like dimensions, one time-like dimension and no degenerate dimension — the numbers of entries $+1$ and $-1$ of $\eta_{71}=\mathrm{diag}(-1,+1,\dots,+1)$ (\texttt{η71}), the same spacetime written in the sign convention of 4+4.nb (space-like $=+1$, time-like $=-1$). (ii) The matrix $\eta=-\eta_{71}=\mathrm{diag}(+1,-1,\dots,-1)$ (\texttt{η17}) is simply the squares of the octonion basis: $e_0^2=+1$ for the unit $e_0=1$ (time) and $e_i^2=-1$ for the seven imaginary units $e_1,\dots,e_7$ (space); more generally it is the real part of the product of two basis units, and it gives the quadratic form of the octonions:
\begin{gather*}
\eta_{AB}=\mathrm{Re}(e_Ae_B),\qquad e_0^2=+1,\qquad e_i^2=-1\ \ (i=1,\dots,7),\\ \mathrm{Re}(x\,x)=x_0^2-x_1^2-\cdots-x_7^2=\eta_{AB}\,x^Ax^B\qquad (x=x^Ae_A\in\mathbb O).
\end{gather*}
$\eta$ gives the signs of the Clifford relation $\Gamma_A\Gamma_B+\Gamma_B\Gamma_A=2\,\eta_{AB}\,1_{16}$ of the real 16×16 Dirac matrices: time squares to $+1_{16}$ and each of the seven space directions to $-1_{16}$, exactly like the octonion units (in the notation $Cl(p,q)$, with $p$ generators squaring to $+1$ and $q$ to $-1$, the algebra $Cl(1,7)\cong\mathbb R(16)$). It also gives the signs of the metric components: $g_{00}>0$ along the time-like direction, negative squares along the space-like directions. Every component computed by 7+1.nb follows from the component matrix of the metric; the declared list fixes for xAct the sign of the determinant of the metric, \texttt{SignDetOfMetric[g71]} $=-1$, which is the same in both sign conventions, since $\det \eta=\det \eta_{71}=-1$. 7+1.nb checks \texttt{SignatureOfMetric[g71]} $=\{7,1,0\}$, \texttt{SignDetOfMetric[g71]} $=-1$, $\mathrm{Re}(e_Ae_B)=\eta_{AB}$ and $\mathrm{Re}(x\,x)=\eta_{AB}x^Ax^B$.

\Needspace{12\baselineskip}\section*{2. Decisions}\addcontentsline{toc}{section}{2. Decisions}
{\footnotesize
\begin{longtable}{>{\raggedright\arraybackslash}p{\dimexpr 0.0500\linewidth-2\tabcolsep\relax}>{\raggedright\arraybackslash}p{\dimexpr 0.4700\linewidth-2\tabcolsep\relax}>{\raggedright\arraybackslash}p{\dimexpr 0.4800\linewidth-2\tabcolsep\relax}}
\toprule
\textbf{\#} & \textbf{decision} & \textbf{reason / equation} \\
\midrule\endfirsthead
\toprule
\textbf{\#} & \textbf{decision} & \textbf{reason / equation} \\
\midrule\endhead
\endfoot
\bottomrule\endlastfoot
1 & append section 25 to 7+1.nb (the auxiliary torus and the world brane) and, for the second request, section 26 (the radion, on $T^4$ and $S^4$, flat and warped); the summary of all checks becomes section 27 & every worked example is computed and checked by the notebook itself, in the same evaluation as everything else \\
2 & write independent Python programs, tools/ref\_torus\_brane.py and tools/ref\_radion.py & their own lattice sums, symbolic integrals, curvature routines and multi-precision arithmetic; no code shared with the notebook \\
3 & back up the evaluated notebook before each appended section and check that the earlier sections are unchanged & tools/verify\_7p1\_nb.py compares every cell and output with the backups (section 6.3) \\
4 & the brane sits at $\vec y=0$ and is not a dynamical field & as in 7+1.nb; all points of the torus are equivalent, so this is a choice of origin \\
5 & the induced metric $h=J^{\mathsf T}gJ$; the brane volume element $\sqrt{|h|}$ & the definition of a submanifold metric \\
6 & brane energy–momentum $T^W_{\mu\nu}=(2/\sqrt{|h|})\,\delta S_W/\delta h^{\mu\nu}$, with $T^0{}_0$ the energy density & the convention of 7+1.nb; it gives $T^W_{\mu\nu}=f_Wh_{\mu\nu}$ for a tension \\
7 & the inhomogeneous metric of section 11.3 is used as the example of a brane that is not totally geodesic & it depends on $x_4$ and is not even in $x_4$, so its Christoffel symbols at the brane do not vanish \\
8 & the brane spinors from the Schur average of the 256 Clifford products & Schur’s lemma: the intertwiner is unique up to a phase \\
9 & the sphere comparison: brane at the north pole of $S^4$, the same volume $V_\perp$ & the same Newton constant, so that only the geometry differs \\
10 & warped branes in five dimensions (RS1), in the conventions $G^A{}_B-\Lambda\delta^A_B=\kappa T^A{}_B$; tensions from the $\delta$-function part of the field equations & the note’s warped formulas are five-dimensional; this form involves no choice of normal \\
11 & the Higgs–radion mixing in the conventions of Giudice, Rattazzi and Wells, translated to the signature $(+,-,-,-)$; the numbers in 60-digit arithmetic & the note follows them; the rate changes, $\sim10^{-31}$, are below double precision \\
12 & in this document the Higgs field is $h$ and the radion $\phi$ & the notation of the Standard Model section; $h_{\mu\nu}$ keeps its indices \\
13 & numbers today with $M_8=10$ TeV, $M_{\rm Pl}=2.435\times10^{18}$ GeV, $m_r=10$ GeV, Planck 2018 & the choices of sections 22–24 \\*
\end{longtable}}
\Needspace{12\baselineskip}\section*{3. Inputs from the earlier sections of 7+1.nb}\addcontentsline{toc}{section}{3. Inputs from the earlier sections of 7+1.nb}
{\footnotesize
\begin{longtable}{>{\raggedright\arraybackslash}p{\dimexpr 0.2600\linewidth-2\tabcolsep\relax}>{\raggedright\arraybackslash}p{\dimexpr 0.2600\linewidth-2\tabcolsep\relax}>{\raggedright\arraybackslash}p{\dimexpr 0.1000\linewidth-2\tabcolsep\relax}>{\raggedright\arraybackslash}p{\dimexpr 0.3800\linewidth-2\tabcolsep\relax}}
\toprule
\textbf{object} & \textbf{Mathematica} & \textbf{section} & \textbf{equation or value} \\
\midrule\endfirsthead
\toprule
\textbf{object} & \textbf{Mathematica} & \textbf{section} & \textbf{equation or value} \\
\midrule\endhead
\endfoot
\bottomrule\endlastfoot
cosmological and inhomogeneous metrics, curvature & \texttt{gCos}, \texttt{gPoly}, \texttt{cvCos}, \texttt{cvPoly} & 11 & $ds^2=n^2dx_0^2-a^2d\vec x^2-b^2d\vec y^2$ \\
Lovelock densities & \texttt{lovelock71.py} & 12 & $L_2$, $L_3$ of the document \\
Standard Model, brane representation, intertwiner & \texttt{g4D}, \texttt{g5D}, \texttt{rhoK}, \texttt{rBr}, \texttt{uInt}, \texttt{rho4567}, \texttt{PL17} & 17 & $\Gamma_AU=UR_A$; $g_*=106.75$; electroweak constants \\
Kaluza–Klein operator, neutrino portal & \texttt{xKK} & 18 & $X=\gamma_5\otimes\vec\rho\cdot\vec q/b_0$; the mass matrix $\mathcal M$ \\
stress tensor, radion, stabilizer, Friedmann equations & \texttt{feStress}, \texttt{uPot} & 20 & $\varepsilon=W/(a^3b^4)$ …; $m_r^2=\kappa K_U/12$; $8\pi G_4=\kappa/V_\perp$ \\
primordial regime, capture, decays, today, far future & \texttt{rIp1}, \texttt{rIq1}, \texttt{rIIr4}, \texttt{rIIIH} & 21–23 & the numbers of sections 7 and 11 of the document \\
natural couplings and their dilution & \texttt{dkLam}, \texttt{dkYeff}, \texttt{dkMuP} & 24 & $\lambda_{\rm eff}=c_PM_8/M_{\rm Pl}^2$, $y_{\rm eff}=c_yM_8/M_{\rm Pl}$ \\
the auxiliary torus and the world brane & \texttt{tb…} & 25 & induced metrics, Christoffel symbols, $U$, brane spinors, smeared sources \\
the radion, the sphere, warped branes, Higgs–radion mixing & \texttt{rd…} & 26 & sections 7.5, 12 and 13 of the document \\*
\end{longtable}}
Before section 25 was appended, the evaluated notebook was copied to the first file below, before section 26 to the second; both are excluded from the repository by .gitignore (size and SHA-256):

\begin{Verbatim}[fontsize=\small,frame=single,framesep=2mm]
backups/7+1_2026-10-28_before_torus_brane.nb
    1940664 bytes, SHA-256 39f64ee81d35c0c9fb94188fd7f8c8e75c96ea198fbf72a5ca6b3c0027426018
backups/7+1_2026-10-28_before_radion.nb
    2031601 bytes, SHA-256 44385a6b1d808f1f95b1d634d29c3960c6251cab96f436b9927522abe0fd8397
\end{Verbatim}
\Needspace{12\baselineskip}\section*{4. Data and their sources}\addcontentsline{toc}{section}{4. Data and their sources}
{\footnotesize
\begin{longtable}{>{\raggedright\arraybackslash}p{\dimexpr 0.2500\linewidth-2\tabcolsep\relax}>{\raggedright\arraybackslash}p{\dimexpr 0.3700\linewidth-2\tabcolsep\relax}>{\raggedright\arraybackslash}p{\dimexpr 0.3800\linewidth-2\tabcolsep\relax}}
\toprule
\textbf{quantity} & \textbf{value used} & \textbf{source} \\
\midrule\endfirsthead
\toprule
\textbf{quantity} & \textbf{value used} & \textbf{source} \\
\midrule\endhead
\endfoot
\bottomrule\endlastfoot
$\hbar c$ & $1.973269804\times10^{-16}$ GeV m & CODATA 2018 (E. Tiesinga et al., Rev. Mod. Phys. 93 (2021) 025010) \\
$G_F$, $m_W$, $m_Z$, $m_H$ & $1.1663788\times10^{-5}$ GeV$^{-2}$, 80.3692, 91.1880, 125.20 GeV & Particle Data Group 2025 (section 17.3) \\
$H_0$, $\Omega_m$, $\Omega_r$ & 67.4 km/s/Mpc, 0.315, $9.14\times10^{-5}$ & Planck 2018 (Astron. Astrophys. 641 (2020) A6) \\
$M_{\rm Pl}$ (reduced), $M_8$, $m_r$ & $2.435\times10^{18}$ GeV, 10 TeV, 10 GeV & as in sections 17–24 \\
RS1 conventions & $2M^3\int\sqrt{-G}R$, $V_{\rm hid}=-V_{\rm vis}=24M^3k$ & L. Randall, R. Sundrum, Phys. Rev. Lett. 83 (1999) 3370 \\
Higgs–radion conventions & $\xi$, $Z^2=1+6\xi\gamma^2(1-6\xi)$ & G. F. Giudice, R. Rattazzi, J. D. Wells, Nucl. Phys. B 595 (2001) 250 \\
$\langle\mu\rangle$ of the warped example & 3 TeV & the author’s note \\*
\end{longtable}}
Software: Wolfram Language 15.0.1 (front end WolframNB.exe; the kernel configuration "Eval71" is a fresh kernel); Python 3.13.17 with numpy 2.5.3, scipy 1.18.1, sympy 1.14.0, mpmath 1.3.0.

\Needspace{12\baselineskip}\section*{5. Derivations}\addcontentsline{toc}{section}{5. Derivations}
\textbf{Induced metric.} With $X^A(x)=(x_0,\dots,x_3,0,0,0,0)$, $J^A{}_\mu=\partial X^A/\partial x^\mu$ is the $8\times4$ matrix $(1_4;0)$, and $h_{\mu\nu}=g_{AB}J^A{}_\mu J^B{}_\nu=g_{\mu\nu}(x,0)$. For the cosmological metric $h=\mathrm{diag}(1,-a^2,-a^2,-a^2)$.

\textbf{Extrinsic curvature.} For $n_{(k)}=b^{-1}\partial_k$ and $g_{\mu k}=0$: $K^{(k)}_{\mu\nu}=\frac12(\mathcal L_{n_{(k)}}g)_{\mu\nu}=\frac1{2b}\partial_kg_{\mu\nu}$; and $\{k\,\mu\nu\}=-\frac12g^{kk}\partial_kg_{\mu\nu}=\frac1{2b^2}\partial_kg_{\mu\nu}$ ($g^{kk}=-b^{-2}$). For the inhomogeneous metric at $x_0=\frac12$, $x_4=0$: $n=\frac{45}{44}$, $\partial_4n=\frac1{10}$, $b=\frac{73}{36}$, so $\{4\,00\}=n\partial_4n/b^2=1458/58619$.

\textbf{Tension.} $\delta\sqrt{|h|}=-\frac12\sqrt{|h|}h_{\mu\nu}\delta h^{\mu\nu}$ gives $T^W_{\mu\nu}=\frac2{\sqrt{|h|}}\frac{\delta(-f_W\sqrt{|h|})}{\delta h^{\mu\nu}}=f_Wh_{\mu\nu}$, hence $\varepsilon=f_W$, $p=-f_W$, $w=-1$.

\textbf{Smeared sources.} The first law $dW=-p_a\,d(a^3b^4)|_b-p_b\,d(a^3b^4)|_a$ gives $\varepsilon=W/(a^3b^4)$, $p_a=-\partial_aW/(3a^2b^4)$, $p_b=-\partial_bW/(4a^3b^3)$; a brane cell has the volume $a^3$, so $W$ does not depend on $b$ and $p_b=0$.

\textbf{Bulk radiation.} For $W=\sqrt{k_a^2/a^2+k_b^2/b^2}$ with momenta isotropic in seven dimensions, the angle between the three- and four-dimensional parts has the weight $\cos^2\theta\sin^3\theta$:
\begin{equation*}
\langle\cos^2\theta\rangle=\frac{\int_0^{\pi/2}\cos^4\theta\sin^3\theta\,d\theta}{\int_0^{\pi/2}\cos^2\theta\sin^3\theta\,d\theta}=\frac{2/35}{2/15}=\frac37,\qquad p_a=\frac13\cdot\frac37\varepsilon=\frac\varepsilon7,\qquad p_b=\frac14\cdot\frac47\varepsilon=\frac\varepsilon7 .
\end{equation*}
\textbf{Images and Kaluza–Klein modes.} Poisson summation: $\sum_{\vec n}(r^2+|\vec n|^2)^{-5/2}=\frac{4\pi^2}{3r}\sum_{\vec m}e^{-2\pi r|\vec m|}$ (lengths in units of $\ell_\perp$), from the four-dimensional Fourier transform $\int d^4y\,e^{-i\vec p\cdot\vec y}(r^2+|\vec y|^2)^{-5/2}=\frac{4\pi^2}{3r}e^{-pr}$; with $8\pi G_4=\kappa/V_\perp$ this gives $\Phi_{\rm N}\to-\frac53G_4M/r$.

\textbf{Radion coupling.} $h=e^{-4\beta}g_E$; a mass at rest: $-M\int e^{-2\beta}dx_0$, linear term $2M\beta=M\phi/(\sqrt3M_{\rm Pl})$; the potential between two masses $-\frac{MM'}{12\pi M_{\rm Pl}^2}\frac{e^{-m_rr}}r=-\frac23\frac{G_4MM'}re^{-m_rr}$.

\textbf{Brane spinors.} $\Gamma_AU=UR_A$ with $R_\mu=\gamma_\mu\otimes1_4$, $R_k=\gamma_5\otimes\rho_k$; then $\Gamma_\mu\Gamma_\nu=U(\gamma_\mu\gamma_\nu\otimes1_4)U^\dagger$, $\Gamma_k\Gamma_l=U(1_4\otimes\rho_k\rho_l)U^\dagger$ ($\gamma_5^2=1_4$), and $\chi=i\Gamma_0\cdots\Gamma_7=U(\gamma_5\otimes\rho_{4567})U^\dagger$ from $R_0R_1R_2R_3=-i\gamma_5\otimes1_4$ and $R_4R_5R_6R_7=1_4\otimes\rho_{4567}$.

\textbf{Kaluza–Klein masses at the brane.} $(\gamma\cdot p\otimes1_4+X)^2=(\gamma\cdot p)^2\otimes1_4+X^2$ because $\gamma_5$ anticommutes with every $\gamma_\mu$, and $X^2=-|\vec q|^2/b_0^2$ from $\rho_k\rho_l+\rho_l\rho_k=-2\delta_{kl}$: $m_{\vec n}^2=m^2+4\pi^2|\vec n|^2/\ell_\perp^2$.

\textbf{Dilution of the portals.} A brane coupling of a bulk field at $\vec y=0$ involves $\phi_{\vec n}(0)=1/\sqrt{V_\perp}$ per bulk field: $\lambda_{\rm eff}=\lambda_P/V_\perp$, $y_{\rm eff}=y/\sqrt{V_\perp}$, $G_{\rm eff}=G_P/V_\perp$; with $V_\perp=M_{\rm Pl}^2/M_8^6$ and natural couplings these are $c_PM_8/M_{\rm Pl}^2$, $c_yM_8/M_{\rm Pl}$, $c_V/M_{\rm Pl}^2$.

\textbf{The normalcy temperature} (ESTIMATE): the emission rate $T^7/M_8^6$ equals the expansion rate $T^2/M_{\rm Pl}$ at $T_*=(M_8^6/M_{\rm Pl})^{1/5}$.

\textbf{The sphere brane.} At the north pole of $S^4$ the induced metric of $ds^2=e^{-4\beta}g_E-e^{2\beta}r_0^2\hat\gamma$ is $e^{-4\beta}g_E$, as on the torus. The rotations $L_{ab}=X_a\partial_b-X_b\partial_a$ vanish at $(0,0,0,0,r)$ iff $a,b\le4$ (6 of 10). The zonal harmonics are normalized by $\int_{S^4}(C_l/C_l(1))^2=V/d_l$ (exact for $l\le6$), so $Z_l(0)=\sqrt{d_l/V}$. The Killing spinors satisfy $\nabla_i\varepsilon_s=s\frac i{2r}e_i{}^a\tilde\gamma_a\varepsilon_s$ exactly (checked after writing the angles as twice half-angles), hence $\not D\varepsilon_s=s\frac{2i}r\varepsilon_s$.

\textbf{Warped branes.} For $ds^2=e^{-2A}\eta-dy^2$: $G^\mu{}_\nu=3A''-6A'^2$; $A=k|y|$ gives $A''=\pm2k\delta$ at the branes and the tensions $\pm6k/\kappa$; with $K^\mu{}_\nu=\frac12g^{\mu\alpha}\partial_yg_{\alpha\nu}=-A'\delta^\mu_\nu$, $[K^\mu{}_\nu]-\delta^\mu_\nu[K]=3[A']$ is the $\delta$ part of $G^\mu{}_\nu$. In eight dimensions with warping along one direction, $G^z{}_z=4A''-10A'^2\ne G^y{}_y=-6A'^2$ unless $A'=0$.

\textbf{Higgs–radion.} For $h=\Omega^2\eta$, $\sqrt{|h|}R(h)=-6\Omega\Box\Omega$ (computed); with $\Omega=1+o_1\phi+\dots$ and $\varphi_H=v+h$ the curvature term has the quadratic part $-3\xi(v^2o_1^2\phi\Box\phi+2vo_1h\Box\phi)$, which after integrations by parts is $3\xi\gamma^2(\partial\phi)^2+6\xi\gamma\,\partial h\cdot\partial\phi$ ($\gamma=vo_1$); the shift and the rotation then give $Z^2$ and $\tan2\theta$ of the document.

\Needspace{12\baselineskip}\section*{6. Verification}\addcontentsline{toc}{section}{6. Verification}
\Needspace{8\baselineskip}\subsection*{6.1 The checks of sections 25 and 26 of 7+1.nb}\addcontentsline{toc}{subsection}{6.1 The checks of sections 25 and 26 of 7+1.nb}
\textbf{Evaluation of 7+1.nb}: Wolfram 15.0.1 for Microsoft Windows (64-bit) (July 2, 2026), a fresh kernel (front-end kernel configuration Eval71); finished Wed 28 Oct 2026 08:57:47; 283.883 s; message cells 0; output cells 554. The checks of section 25 (1–14 the torus, 15–16 torus and brane, 17–27 the world brane):

{\footnotesize
\begin{longtable}{>{\raggedright\arraybackslash}p{\dimexpr 0.0550\linewidth-2\tabcolsep\relax}>{\raggedright\arraybackslash}p{\dimexpr 0.7450\linewidth-2\tabcolsep\relax}>{\raggedright\arraybackslash}p{\dimexpr 0.0800\linewidth-2\tabcolsep\relax}>{\raggedright\arraybackslash}p{\dimexpr 0.1200\linewidth-2\tabcolsep\relax}}
\toprule
\textbf{\#} & \textbf{check} & \textbf{result} & \textbf{section} \\
\midrule\endfirsthead
\toprule
\textbf{\#} & \textbf{check} & \textbf{result} & \textbf{section} \\
\midrule\endhead
\endfoot
\bottomrule\endlastfoot
1 & \texttt{auxiliary torus: for the cosmological metric det(-g\_\allowbreak{}kl) == b\textasciicircum{}8 (k, l = 4..7), and the proper volume of the fundamental domain [0, Lc)\textasciicircum{}4, Integrate[Sqrt[det(-g\_\allowbreak{}kl)], \{y, 0, Lc\}\textasciicircum{}4], is b\textasciicircum{}4 Lc\textasciicircum{}4 = V\_\allowbreak{}perp} & True & 25.1 \\
2 & \texttt{only the product b Lc of the auxiliary torus is physical: y = lambda y' maps the torus of period Lc and size b to the one of period Lc/\allowbreak{}lambda and size lambda b; the line element b\textasciicircum{}2 dy.dy, the proper period b Lc and V\_\allowbreak{}perp = b\textasciicircum{}4 Lc\textasciicircum{}4 are unchanged} & True & 25.1 \\
3 & \texttt{symmetries of the auxiliary torus that fix y = 0 (and the brane): the 384 = 2\textasciicircum{}4 4! signed permutation matrices are distinct orthogonal integer matrices (they map Z\textasciicircum{}4 onto itself and preserve |n|\textasciicircum{}2); the shells |n|\textasciicircum{}2 = 1, 2, 3, 4, 9 split into orbits of sizes \{8\}, \{24\}, \{32\}, \{8, 16\}, \{8, 96\}} & True & 25.1 \\
4 & \texttt{Fourier modes of the auxiliary torus: phi\_\allowbreak{}n = Exp[I q\_\allowbreak{}n.y]/\allowbreak{}Sqrt[V\_\allowbreak{}perp] with q\_\allowbreak{}n = 2 Pi n/\allowbreak{}Lc (n in Z\textasciicircum{}4) are orthonormal for the proper volume element b\textasciicircum{}4 d\textasciicircum{}4y (exact integrals: n = m = (1, 0, -2, 3); n = (1, 0, -2, 3), m = (1, 1, -2, 3); n = 0, m = (0, 2, 0, 0))} & True & 25.2 \\
5 & \texttt{Parseval on the auxiliary torus: f(y) = Prod\_\allowbreak{}k cos\textasciicircum{}2(Pi y\_\allowbreak{}k/\allowbreak{}Lc) = Prod\_\allowbreak{}k (1/\allowbreak{}2 + (1/\allowbreak{}4) e\textasciicircum{}(2 Pi I y\_\allowbreak{}k/\allowbreak{}Lc) + (1/\allowbreak{}4) e\textasciicircum{}(-2 Pi I y\_\allowbreak{}k/\allowbreak{}Lc)) has the 81 non-zero coefficients c\_\allowbreak{}n = Sqrt[V\_\allowbreak{}perp] Prod\_\allowbreak{}k w(n\_\allowbreak{}k), w(0) = 1/\allowbreak{}2, w(+-1) = 1/\allowbreak{}4 (three of them checked by integration), and Sum\_\allowbreak{}n |c\_\allowbreak{}n|\textasciicircum{}2 == Int b\textasciicircum{}4 d\textasciicircum{}4y f\textasciicircum{}2 == (3/\allowbreak{}8)\textasciicircum{}4 V\_\allowbreak{}perp} & True & 25.2 \\
6 & \texttt{the Laplacian of the auxiliary torus: -(1/\allowbreak{}b\textasciicircum{}2) Sum\_\allowbreak{}k d\textasciicircum{}2 phi\_\allowbreak{}n/\allowbreak{}dy\_\allowbreak{}k\textasciicircum{}2 == (|q\_\allowbreak{}n|\textasciicircum{}2/\allowbreak{}b\textasciicircum{}2) phi\_\allowbreak{}n (n = (2, -1, 0, 3))} & True & 25.2 \\
7 & \texttt{scalar Kaluza-Klein masses: for the static metric diag(1, -1, -1, -1, -b0\textasciicircum{}2, ..., -b0\textasciicircum{}2) the Klein-Gordon operator Box + m\textasciicircum{}2 (Box f = |g|\textasciicircum{}(-1/\allowbreak{}2) d\_\allowbreak{}A(|g|\textasciicircum{}(1/\allowbreak{}2) g\textasciicircum{}AB d\_\allowbreak{}B f)) on Exp[-I E x0 + I k.x + I q.y] gives (-E\textasciicircum{}2 + |k|\textasciicircum{}2 + |q|\textasciicircum{}2/\allowbreak{}b0\textasciicircum{}2 + m\textasciicircum{}2) times the wave: in ordinary spacetime a particle of mass Sqrt[m\textasciicircum{}2 + |q|\textasciicircum{}2/\allowbreak{}b0\textasciicircum{}2]} & True & 25.2 \\
8 & \texttt{Jacobi's four-square theorem: the number of n in Z\textasciicircum{}4 with |n|\textasciicircum{}2 = k is r\_\allowbreak{}4(k) = 8 sigma(k) - 32 sigma(k/\allowbreak{}4) for k = 1..100 (direct count over \{-10, ..., 10\}\textasciicircum{}4)} & True & 25.3 \\
9 & \texttt{Weyl's law: the number N(R) of n in Z\textasciicircum{}4 with |n| <= R approaches the volume Pi\textasciicircum{}2 R\textasciicircum{}4/\allowbreak{}2 of the four-ball: N(10) = 49689, N(100) = 493490641, and |N(R)/\allowbreak{}(Pi\textasciicircum{}2 R\textasciicircum{}4/\allowbreak{}2) - 1| < 10\textasciicircum{}-2, 10\textasciicircum{}-3, 3 10\textasciicircum{}-5, 5 10\textasciicircum{}-6, 10\textasciicircum{}-6 for R = 10, 30, 100, 300, 1000} & True & 25.3 \\
10 & \texttt{the auxiliary torus today (M8 = 10 TeV, M\_\allowbreak{}Pl = 2.435 10\textasciicircum{}18 GeV): V\_\allowbreak{}perp = M\_\allowbreak{}Pl\textasciicircum{}2/\allowbreak{}M8\textasciicircum{}6 = 5.929 10\textasciicircum{}12 GeV\textasciicircum{}-4, l\_\allowbreak{}perp = V\_\allowbreak{}perp\textasciicircum{}(1/\allowbreak{}4) = b0 Lc = 1560.4 GeV\textasciicircum{}-1 = 3.079 10\textasciicircum{}-13 m, lowest Kaluza-Klein mass 2 Pi/\allowbreak{}l\_\allowbreak{}perp = 4.027 MeV; N(mu) = V\_\allowbreak{}perp mu\textasciicircum{}4/\allowbreak{}(32 Pi\textasciicircum{}2) = 1.17 10\textasciicircum{}13 levels below m\_\allowbreak{}r/\allowbreak{}2 = 5 GeV (the lattice count of section 22.2 agrees to 10\textasciicircum{}-5) and N(M8) = M\_\allowbreak{}Pl\textasciicircum{}2/\allowbreak{}(32 Pi\textasciicircum{}2 M8\textasciicircum{}2)} & True & 25.3 \\
11 & \texttt{Newton's law in eight dimensions: Omega\_\allowbreak{}6 = 2 Pi\textasciicircum{}(7/\allowbreak{}2)/\allowbreak{}Gamma(7/\allowbreak{}2) = 16 Pi\textasciicircum{}3/\allowbreak{}15, r\textasciicircum{}-5 is harmonic away from r = 0 in seven space dimensions ((1/\allowbreak{}r\textasciicircum{}6) d/\allowbreak{}dr (r\textasciicircum{}6 d/\allowbreak{}dr r\textasciicircum{}-5) = 0), and grad\textasciicircum{}2 Phi = 8 Pi G8 ((D - 3)/\allowbreak{}(D - 2)) M delta\textasciicircum{}7 (D = 8) gives Phi = -(5/\allowbreak{}6) kappa M/\allowbreak{}(5 Omega\_\allowbreak{}6 r\textasciicircum{}5) = -5 kappa M/\allowbreak{}(32 Pi\textasciicircum{}3 r\textasciicircum{}5); the same formula in D = 4 gives -G M/\allowbreak{}r} & True & 25.4 \\
12 & \texttt{images = Kaluza-Klein modes (Poisson summation; lengths in units of l\_\allowbreak{}perp): Sum\_\allowbreak{}n (r\textasciicircum{}2 + |n|\textasciicircum{}2)\textasciicircum{}(-5/\allowbreak{}2) == (4 Pi\textasciicircum{}2/\allowbreak{}(3 r)) Sum\_\allowbreak{}m Exp[-2 Pi r |m|] to 10\textasciicircum{}-8 at r = 1/\allowbreak{}4, 1/\allowbreak{}2, 1, 2 (left: shells with r\_\allowbreak{}4(k) to k = 1.5 10\textasciicircum{}6 and the continuum tail Pi\textasciicircum{}2 Int k (r\textasciicircum{}2 + k)\textasciicircum{}(-5/\allowbreak{}2) dk; right: shells to k = 20000)} & True & 25.4 \\
13 & \texttt{near and far field: Z(5/\allowbreak{}2) = Sum\_\allowbreak{}(n != 0) |n|\textasciicircum{}-5 = Sum\_\allowbreak{}k r\_\allowbreak{}4(k) k\textasciicircum{}(-5/\allowbreak{}2) = 7 zeta(5/\allowbreak{}2) zeta(3/\allowbreak{}2) = 24.531 (shell sum with tail, to 10\textasciicircum{}-8); near the source r\textasciicircum{}5 S(r) - 1 -> r\textasciicircum{}5 Z(5/\allowbreak{}2) (the ratio lies between 0.98 and 1 at r = 1/\allowbreak{}10 and 1/\allowbreak{}20, closer to 1 at 1/\allowbreak{}20): the eight-dimensional law; far from it 3 r S/\allowbreak{}(4 Pi\textasciicircum{}2) - 1 -> 8 Exp[-2 Pi r], the 8 lightest Kaluza-Klein modes (ratio 1 to 2 10\textasciicircum{}-2 at r = 2 and 3)} & True & 25.4 \\
14 & \texttt{graviton + radion: far from the source the images add up to (5 kappa M/\allowbreak{}(32 Pi\textasciicircum{}3)) (4 Pi\textasciicircum{}2/\allowbreak{}(3 V\_\allowbreak{}perp r)) == (5/\allowbreak{}3) G4 M/\allowbreak{}r with G4 = kappa/\allowbreak{}(8 Pi V\_\allowbreak{}perp); the radion of section 20.3, coupled to matter by (phi/\allowbreak{}(Sqrt[3] M\_\allowbreak{}Pl)) T, gives the potential -(1/\allowbreak{}(Sqrt[3] M\_\allowbreak{}Pl))\textasciicircum{}2 M M' Exp[-m\_\allowbreak{}r r]/\allowbreak{}(4 Pi r), and (1/\allowbreak{}(Sqrt[3] M\_\allowbreak{}Pl))\textasciicircum{}2/\allowbreak{}(4 Pi) == (2/\allowbreak{}3) G4 (8 Pi G4 = 1/\allowbreak{}M\_\allowbreak{}Pl\textasciicircum{}2): graviton 1 + massless radion 2/\allowbreak{}3 = 5/\allowbreak{}3} & True & 25.4 \\
15 & \texttt{Kaluza-Klein momentum is conserved in the bulk and not at the brane: Int b\textasciicircum{}4 d\textasciicircum{}4y phi\_\allowbreak{}n phi\_\allowbreak{}m phi\_\allowbreak{}l\textasciicircum{}* == 1/\allowbreak{}Sqrt[V\_\allowbreak{}perp] for l = n + m and 0 otherwise (exact: n = (1, 0, 0, 2), m = (0, 1, -1, 0), l = (1, 1, -1, 2) and (1, 1, -1, 1)); at y = 0, phi\_\allowbreak{}n phi\_\allowbreak{}m phi\_\allowbreak{}l\textasciicircum{}* == V\_\allowbreak{}perp\textasciicircum{}(-3/\allowbreak{}2) for every n, m, l} & True & 25.5 \\
16 & \texttt{the delta function of the torus: Int\_\allowbreak{}(-Lc/\allowbreak{}2)\textasciicircum{}(Lc/\allowbreak{}2) Exp[-I q\_\allowbreak{}n y] delta(y) dy == 1 for every integer n, so delta\textasciicircum{}4(y) = Lc\textasciicircum{}-4 Sum\_\allowbreak{}n Exp[I q\_\allowbreak{}n.y] (every mode with the weight 1/\allowbreak{}Lc\textasciicircum{}4); a brane source rho4 delta\textasciicircum{}4(y)/\allowbreak{}b\textasciicircum{}4 has Int b\textasciicircum{}4 d\textasciicircum{}4y (rho4 delta\textasciicircum{}4(y)/\allowbreak{}b\textasciicircum{}4) == rho4, i.e. the average rho4/\allowbreak{}V\_\allowbreak{}perp over the torus} & True & 25.5 \\
17 & \texttt{the world brane W = \{x4 = x5 = x6 = x7 = 0\}: the induced metric h = J\textasciicircum{}T.g.J of the embedding X = (x0, x1, x2, x3, 0, 0, 0, 0) is, for the cosmological metric of section 11.2, the Friedmann-Robertson-Walker metric diag(1, -a\textasciicircum{}2, -a\textasciicircum{}2, -a\textasciicircum{}2) (independent of b); for the inhomogeneous metric of section 11.3 it is diag(n\textasciicircum{}2, -a\textasciicircum{}2, -a\textasciicircum{}2, -a\textasciicircum{}2) at x4 = 0 with n = 1 + x0\textasciicircum{}2/\allowbreak{}11, a = 1 + x0/\allowbreak{}3} & True & 25.6 \\
18 & \texttt{brane integrals in eight-dimensional form: at the brane Sqrt|det g| == Sqrt|det h| b\textasciicircum{}4 with b\textasciicircum{}4 = Sqrt[det(-g\_\allowbreak{}kl)] (cosmological metric; inhomogeneous metric exactly at x0 = 1/\allowbreak{}2), so Int\_\allowbreak{}W d\textasciicircum{}4x Sqrt|h| L == Int d\textasciicircum{}8X Sqrt|g| (delta\textasciicircum{}4(y)/\allowbreak{}b\textasciicircum{}4) L; test: Int d\textasciicircum{}4y b\textasciicircum{}4 (delta\textasciicircum{}4(y)/\allowbreak{}b\textasciicircum{}4) F(y) == F(0) = 1 for F = Exp[x4 + 2 x5] Cos[x6 - x7]} & True & 25.6 \\
19 & \texttt{how the brane sits in the bulk: the Christoffel symbols \{k mu nu\} at the brane (k = 4..7, mu, nu = 0..3; = (1/\allowbreak{}(2 b\textasciicircum{}2)) d\_\allowbreak{}k g\_\allowbreak{}mu nu for these metrics, b times the extrinsic curvature) all vanish for the cosmological metric: the brane is totally geodesic, and a free particle moving along it stays on it. For the inhomogeneous metric at x0 = 1/\allowbreak{}2 the only non-zero one is \{4 00\} = n (d\_\allowbreak{}4 n)/\allowbreak{}b\textasciicircum{}2 = 1458/\allowbreak{}58619: a free particle at rest on the brane is accelerated off it} & True & 25.6 \\
20 & \texttt{explicit forms: every Gamma17[A] is a signed permutation matrix (one entry +-1 in each row and column), and 2 Sqrt[2] U has only the entries 0, +-1, +-I, eight non-zero ones in each row and in each column} & True & 25.7 \\
21 & \texttt{brane spinor (x) transverse spinor: Gamma17[mu].Gamma17[nu] == U.(gamma\_\allowbreak{}mu.gamma\_\allowbreak{}nu (x) 1\_\allowbreak{}4).U\textasciicircum{}dag, Gamma17[k].Gamma17[l] == U.(1\_\allowbreak{}4 (x) rho\_\allowbreak{}k.rho\_\allowbreak{}l).U\textasciicircum{}dag and Gamma17[mu].Gamma17[k] == U.(gamma\_\allowbreak{}mu.gamma5 (x) rho\_\allowbreak{}k).U\textasciicircum{}dag (mu, nu = 0..3, k, l = 4..7): the Lorentz group of the brane acts on the brane spinor index, the rotations about the brane on the transverse index, and the generators that mix them are broken by the brane} & True & 25.7 \\
22 & \texttt{the transverse index carries Spin(4) = SU(2) x SU(2): rho4567 == -gamma5 (4 x 4), it commutes with every rho\_\allowbreak{}k.rho\_\allowbreak{}l and anticommutes with every rho\_\allowbreak{}k; for each even permutation (k, l, m, n) of (4, 5, 6, 7) the combination rho\_\allowbreak{}k.rho\_\allowbreak{}l + s rho\_\allowbreak{}m.rho\_\allowbreak{}n vanishes on the eigenspace rho4567 = s, and rho\_\allowbreak{}k.rho\_\allowbreak{}l - s rho\_\allowbreak{}m.rho\_\allowbreak{}n acts there with rank 2: the four values of a form (2, 1) + (1, 2)} & True & 25.7 \\
23 & \texttt{chirality at the brane: PL17 == U.(P\_\allowbreak{}L (x) P\_\allowbreak{}+ + P\_\allowbreak{}R (x) P\_\allowbreak{}-).U\textasciicircum{}dag with P\_\allowbreak{}L,R = (1 -+ gamma5)/\allowbreak{}2 and P\_\allowbreak{}+- = (1 +- rho4567)/\allowbreak{}2: the left-handed half psi\_\allowbreak{}L (chi17 = -1, 8 components) is two brane-left-handed Weyl spinors with rho4567 = +1 and two brane-right-handed ones with rho4567 = -1 (each part of rank 4)} & True & 25.7 \\
24 & \texttt{an explicit spinor: psi = U.(u (x) e\_\allowbreak{}1) with the brane spinor u = (1, 2, 0, -1): psibar psi == ubar u == 4 and psibar Gamma\textasciicircum{}mu psi == ubar gamma\textasciicircum{}mu u (mu = 0..3)} & True & 25.7 \\
25 & \texttt{brane tension: for S\_\allowbreak{}W = -Int d\textasciicircum{}4x Sqrt|h| f\_\allowbreak{}W and h = diag(1/\allowbreak{}u0, -1/\allowbreak{}u1, -1/\allowbreak{}u2, -1/\allowbreak{}u3) (h\textasciicircum{}00 = u0, h\textasciicircum{}ii = -u\_\allowbreak{}i), T\_\allowbreak{}mu nu = (2/\allowbreak{}Sqrt|h|) dS\_\allowbreak{}W/\allowbreak{}dh\textasciicircum{}mu nu == f\_\allowbreak{}W h\_\allowbreak{}mu nu, T\textasciicircum{}0\_\allowbreak{}0 == f\_\allowbreak{}W (energy density) and T\textasciicircum{}1\_\allowbreak{}1 == f\_\allowbreak{}W (pressure -f\_\allowbreak{}W: w = -1 along the brane)} & True & 25.8 \\
26 & \texttt{sources smeared over the torus (W(a, b) the comoving energy, section 20.2): the brane tension W = a\textasciicircum{}3 f\_\allowbreak{}W/\allowbreak{}Lc\textasciicircum{}4 gives eps = f\_\allowbreak{}W/\allowbreak{}V\_\allowbreak{}perp, pa = -eps, pb = 0; brane radiation W = a\textasciicircum{}3 r0 a\textasciicircum{}-4/\allowbreak{}Lc\textasciicircum{}4 gives pa = eps/\allowbreak{}3, pb = 0; radiation of the bulk, isotropic in the seven space directions, has <cos\textasciicircum{}2 theta> = 3/\allowbreak{}7 and <sin\textasciicircum{}2 theta> = 4/\allowbreak{}7 for the weight cos\textasciicircum{}2 theta sin\textasciicircum{}3 theta of the momentum sphere, so pa = (1/\allowbreak{}3)(3/\allowbreak{}7) eps = eps/\allowbreak{}7 and pb = (1/\allowbreak{}4)(4/\allowbreak{}7) eps = eps/\allowbreak{}7} & True & 25.8 \\
27 & \texttt{the brane in four-dimensional physics: (i) the tension is vacuum energy, kappa f\_\allowbreak{}W/\allowbreak{}V\_\allowbreak{}perp == f\_\allowbreak{}W/\allowbreak{}M\_\allowbreak{}Pl\textasciicircum{}2 (V\_\allowbreak{}perp = kappa M\_\allowbreak{}Pl\textasciicircum{}2), and f\_\allowbreak{}W = M8\textasciicircum{}4 is between 3 10\textasciicircum{}62 and 5 10\textasciicircum{}62 times rho\_\allowbreak{}Lambda = 3 H0\textasciicircum{}2 M\_\allowbreak{}Pl\textasciicircum{}2 Omega\_\allowbreak{}Lambda (section 24.5); (ii) normalcy temperature: (T\textasciicircum{}7/\allowbreak{}M8\textasciicircum{}6)/\allowbreak{}(T\textasciicircum{}2/\allowbreak{}M\_\allowbreak{}Pl) == 1 at T* = (M8\textasciicircum{}6/\allowbreak{}M\_\allowbreak{}Pl)\textasciicircum{}(1/\allowbreak{}5) = 13.3 GeV, and V\_\allowbreak{}perp T\textasciicircum{}4 T\textasciicircum{}3/\allowbreak{}M\_\allowbreak{}Pl\textasciicircum{}2 == T\textasciicircum{}7/\allowbreak{}M8\textasciicircum{}6; (iii) thin brane: (1/\allowbreak{}y\textasciicircum{}3) d/\allowbreak{}dy (y\textasciicircum{}3 d/\allowbreak{}dy y\textasciicircum{}-2) == 0 and 2 Pi\textasciicircum{}2 y\textasciicircum{}3 d/\allowbreak{}dy y\textasciicircum{}-2 == -4 Pi\textasciicircum{}2 (the Green's function -1/\allowbreak{}(4 Pi\textasciicircum{}2 |y|\textasciicircum{}2) of the four transverse directions), and Sqrt[kappa f\_\allowbreak{}W] == 1/\allowbreak{}M8 for f\_\allowbreak{}W = M8\textasciicircum{}4} & True & 25.8 \\*
\end{longtable}}
The checks of section 26 (the radion; sections 7.5, 12 and 13 of the document):

{\footnotesize
\begin{longtable}{>{\raggedright\arraybackslash}p{\dimexpr 0.0550\linewidth-2\tabcolsep\relax}>{\raggedright\arraybackslash}p{\dimexpr 0.7450\linewidth-2\tabcolsep\relax}>{\raggedright\arraybackslash}p{\dimexpr 0.0800\linewidth-2\tabcolsep\relax}>{\raggedright\arraybackslash}p{\dimexpr 0.1200\linewidth-2\tabcolsep\relax}}
\toprule
\textbf{\#} & \textbf{check} & \textbf{result} & \textbf{section} \\
\midrule\endfirsthead
\toprule
\textbf{\#} & \textbf{check} & \textbf{result} & \textbf{section} \\
\midrule\endhead
\endfoot
\bottomrule\endlastfoot
1 & \texttt{section 26 engine for diagonal metrics: rdScal and rdEin reproduce curvature[] of section 11 for the cosmological metric (R and every G\textasciicircum{}A\_\allowbreak{}A)} & True & 26.1 \\
2 & \texttt{the 4-sphere of radius r0: Ricci\_\allowbreak{}ij = 3 gamma\_\allowbreak{}ij (gamma the unit-sphere metric), Ricci\_\allowbreak{}ij = 0 for i != j, so R(S\textasciicircum{}4) = 12/\allowbreak{}r0\textasciicircum{}2 (in the signature of 7+1.nb the space-like block gives R = -12/\allowbreak{}r0\textasciicircum{}2, also with curvature[] of section 11); the volume of the unit S\textasciicircum{}4 is 8 Pi\textasciicircum{}2/\allowbreak{}3} & True & 26.1 \\
3 & \texttt{radion on T\textasciicircum{}4: for ds\textasciicircum{}2 = e\textasciicircum{}(-4 beta) (N\textasciicircum{}2 dx0\textasciicircum{}2 - a\textasciicircum{}2 dx.dx) - e\textasciicircum{}(2 beta) b0\textasciicircum{}2 dy.dy the reduced action -(1/\allowbreak{}(2 kappa)) Int(R + 2 Lambda) - Int U equals -(M\_\allowbreak{}Pl\textasciicircum{}2/\allowbreak{}2) R\_\allowbreak{}E + 6 M\_\allowbreak{}Pl\textasciicircum{}2 (d beta)\textasciicircum{}2 - e\textasciicircum{}(-4 beta)(M\_\allowbreak{}Pl\textasciicircum{}2 Lambda + V0 U) up to a total derivative (Euler-Lagrange expressions for N, a, beta vanish), M\_\allowbreak{}Pl\textasciicircum{}2 = V0/\allowbreak{}kappa, V0 = b0\textasciicircum{}4 Lc\textasciicircum{}4} & True & 26.1 \\
4 & \texttt{radion on S\textasciicircum{}4: the same reduction for ds\textasciicircum{}2 = e\textasciicircum{}(-4 beta) g\_\allowbreak{}E - e\textasciicircum{}(2 beta) r0\textasciicircum{}2 gamma (unit-sphere gamma) holds with the additional potential -(M\_\allowbreak{}Pl\textasciicircum{}2/\allowbreak{}2) e\textasciicircum{}(-4 beta) 12/\allowbreak{}(r0\textasciicircum{}2 e\textasciicircum{}(2 beta)) = -6 M\_\allowbreak{}Pl\textasciicircum{}2 e\textasciicircum{}(-6 beta)/\allowbreak{}r0\textasciicircum{}2 from the curvature of the sphere, V0 = (8 Pi\textasciicircum{}2/\allowbreak{}3) r0\textasciicircum{}4; R is independent of the angles} & True & 26.1 \\
5 & \texttt{n extra dimensions (torus T\textasciicircum{}n, D = 4 + n, n = 1..4, Einstein frame e\textasciicircum{}(-n B) g\_\allowbreak{}E): the kinetic term is (M\_\allowbreak{}Pl\textasciicircum{}2/\allowbreak{}2)(n(n+2)/\allowbreak{}2)(dB)\textasciicircum{}2; with the note's B = beta\_\allowbreak{}n phi, beta\_\allowbreak{}n = Sqrt[2/\allowbreak{}(n(n+2))], it is (M\_\allowbreak{}Pl\textasciicircum{}2/\allowbreak{}2)(d phi)\textasciicircum{}2, not (M\_\allowbreak{}Pl\textasciicircum{}2/\allowbreak{}2)(1/\allowbreak{}2)(d phi)\textasciicircum{}2; the form with 1/\allowbreak{}2 needs beta\_\allowbreak{}n = 1/\allowbreak{}Sqrt[n(n+2)]} & True & 26.1 \\
6 & \texttt{the note's eight-dimensional values alpha = -1/\allowbreak{}Sqrt[3], beta = 1/\allowbreak{}Sqrt[12]: 2 alpha + 4 beta == 0 (Einstein frame), 4 beta == 2/\allowbreak{}Sqrt[3] (V = V0 e\textasciicircum{}(4 beta phi)), and Sqrt[12] beta == 1, so M\_\allowbreak{}Pl phi\_\allowbreak{}note = Sqrt[12] M\_\allowbreak{}Pl ln(b/\allowbreak{}b0) is the canonical radion of section 20.3} & True & 26.1 \\
7 & \texttt{the cosmological metric with the sphere of radius b(x0) (ds\textasciicircum{}2 = dx0\textasciicircum{}2 - a\textasciicircum{}2 dx.dx - b\textasciicircum{}2 dOmega\_\allowbreak{}4\textasciicircum{}2): G\textasciicircum{}A\_\allowbreak{}A = (torus value) + (6, 6, 6, 6, 3, 3, 3, 3)/\allowbreak{}b\textasciicircum{}2 and R = (torus value) - 12/\allowbreak{}b\textasciicircum{}2} & True & 26.2 \\
8 & \texttt{radion equation on S\textasciicircum{}4 (Einstein gravity; 11 and 44 equations with the 00 constraint): Hb' + 3 Ha Hb + 4 Hb\textasciicircum{}2 == kappa (eps/\allowbreak{}6 - pa/\allowbreak{}2 + pb/\allowbreak{}3) + Lambda/\allowbreak{}3 - 3/\allowbreak{}b\textasciicircum{}2} & True & 26.2 \\
9 & \texttt{radion on T\textasciicircum{}4 with U = U0 + U1 beta + (K\_\allowbreak{}U/\allowbreak{}2) beta\textasciicircum{}2: a flat vacuum at b = b0 needs Lambda + kappa U0 = 0 and U1 = 0, and m\_\allowbreak{}r\textasciicircum{}2 = V''(beta)/\allowbreak{}(12 M\_\allowbreak{}Pl\textasciicircum{}2) = kappa K\_\allowbreak{}U/\allowbreak{}12 (section 20.3)} & True & 26.2 \\
10 & \texttt{radion on S\textasciicircum{}4 with U: Lambda + kappa U0 == 6/\allowbreak{}r0\textasciicircum{}2, kappa U1 == -12/\allowbreak{}r0\textasciicircum{}2 and m\_\allowbreak{}r\textasciicircum{}2 == kappa K\_\allowbreak{}U/\allowbreak{}12 - 2/\allowbreak{}r0\textasciicircum{}2 from V(beta); the eight-dimensional radion equation gives the same static solution and the same mass} & True & 26.2 \\
11 & \texttt{S\textasciicircum{}4 with Lambda only: dS4 x S\textasciicircum{}4 with r0\textasciicircum{}2 = 9/\allowbreak{}Lambda, H\textasciicircum{}2 = 1/\allowbreak{}r0\textasciicircum{}2 = Lambda/\allowbreak{}9 solves the 00, 11 and 44 equations; it is a MAXIMUM of V: m\textasciicircum{}2 = -6 H\textasciicircum{}2 from V(beta), and m\textasciicircum{}2 = -6/\allowbreak{}r0\textasciicircum{}2 = -6 H\textasciicircum{}2 from the eight-dimensional radion equation} & True & 26.2 \\
12 & \texttt{S\textasciicircum{}4 with a 4-form flux (illustration; W = a\textasciicircum{}3 b\textasciicircum{}4 q\textasciicircum{}2/\allowbreak{}(2 b\textasciicircum{}8) gives eps = rho\_\allowbreak{}F, pa = -rho\_\allowbreak{}F, pb = +rho\_\allowbreak{}F): Minkowski x S\textasciicircum{}4 needs Lambda = 9/\allowbreak{}(2 r0\textasciicircum{}2) and kappa rho\_\allowbreak{}F = 3/\allowbreak{}(2 r0\textasciicircum{}2); m\textasciicircum{}2 = 6/\allowbreak{}r0\textasciicircum{}2 from V(beta) and from the eight-dimensional radion equation} & True & 26.2 \\
13 & \texttt{the note's Goldberger-Wise potential in eight dimensions: for a static profile Phi(y) (here Phi(x4, x5)) Sqrt|G| (1/\allowbreak{}2 G\textasciicircum{}AB dPhi dPhi - 1/\allowbreak{}2 m\textasciicircum{}2 Phi\textasciicircum{}2) == -N a\textasciicircum{}3 e\textasciicircum{}(-4 beta) Sqrt[gamma] (1/\allowbreak{}2 e\textasciicircum{}(-2 beta) |grad Phi|\textasciicircum{}2\_\allowbreak{}gamma + 1/\allowbreak{}2 m\textasciicircum{}2 Phi\textasciicircum{}2): the Einstein-frame factor is e\textasciicircum{}(-4 beta) (e\textasciicircum{}(-4 beta phi) in the note's notation), not e\textasciicircum{}(+4 beta phi)} & True & 26.2 \\
14 & \texttt{the radion couples to a mass M on the brane as (phi/\allowbreak{}(Sqrt[3] M\_\allowbreak{}Pl)) M: -M e\textasciicircum{}(-n beta/\allowbreak{}2) = -M + (n/\allowbreak{}2) M beta + ..., with beta = phi/\allowbreak{}(M\_\allowbreak{}Pl Sqrt[n (n + 2)/\allowbreak{}2]) the coupling is Sqrt[n/\allowbreak{}(2 (n + 2))]/\allowbreak{}M\_\allowbreak{}Pl = 1/\allowbreak{}(Sqrt[3] M\_\allowbreak{}Pl) for n = 4, and the fifth-force strength 2 (coupling M\_\allowbreak{}Pl)\textasciicircum{}2 = n/\allowbreak{}(n + 2) = 2/\allowbreak{}3} & True & 26.3 \\
15 & \texttt{Eot-Wash (|alpha| = 1 excluded for lambda > 38.6 micrometres): a radion of strength 2/\allowbreak{}3 must be heavier than about hbar c/\allowbreak{}38.6 micrometres = 5.1 meV; the 10 GeV radion has the range 2.0 10\textasciicircum{}-17 m; Higgs-radion mixing in 7+1 is governed by gamma = v/\allowbreak{}(Sqrt[3] M\_\allowbreak{}Pl) = 5.8 10\textasciicircum{}-17} & True & 26.3 \\
16 & \texttt{zonal harmonics on S\textasciicircum{}4 (l = 0..6, exact): (1/\allowbreak{}sin\textasciicircum{}3 x4) d/\allowbreak{}dx4 (sin\textasciicircum{}3 x4 d/\allowbreak{}dx4 C\_\allowbreak{}l\textasciicircum{}(3/\allowbreak{}2)(cos x4)) == -l (l + 3) C\_\allowbreak{}l\textasciicircum{}(3/\allowbreak{}2), C\_\allowbreak{}l\textasciicircum{}(3/\allowbreak{}2)(1) == (l + 1)(l + 2)/\allowbreak{}2, and Int over S\textasciicircum{}4 of (C\_\allowbreak{}l/\allowbreak{}C\_\allowbreak{}l(1))\textasciicircum{}2 == V/\allowbreak{}d\_\allowbreak{}l (unit sphere), so the normalized zonal harmonic is Sqrt[d\_\allowbreak{}l/\allowbreak{}V] at the north pole; d\_\allowbreak{}l == Binomial[l + 4, 4] - Binomial[l + 2, 4] for l = 0..20} & True & 26.4 \\
17 & \texttt{Weyl's law on S\textasciicircum{}4: the number of levels with l(l + 3) <= (mu r)\textasciicircum{}2 approaches (mu r)\textasciicircum{}4/\allowbreak{}12 = V mu\textasciicircum{}4/\allowbreak{}(32 Pi\textasciicircum{}2) (V = 8 Pi\textasciicircum{}2 r\textasciicircum{}4/\allowbreak{}3): relative deviation -1.0 10\textasciicircum{}-2, -1.0 10\textasciicircum{}-4, -1.0 10\textasciicircum{}-6 at mu r = 10, 100, 1000 (like -1/\allowbreak{}(mu r)\textasciicircum{}2)} & True & 26.4 \\
18 & \texttt{the brane at the north pole of S\textasciicircum{}4: the embedding (r0 sin x4 sin x5 sin x6 cos x7, ..., r0 cos x4) of S\textasciicircum{}4 in R\textasciicircum{}5 induces r0\textasciicircum{}2 gamma; the north pole x4 = 0 is (0, 0, 0, 0, r0); of the 10 rotations L\_\allowbreak{}ab of R\textasciicircum{}5 (the Killing vectors of S\textasciicircum{}4, SO(5)) the 6 with a, b <= 4 vanish there (SO(4), kept by the brane) and the 4 L\_\allowbreak{}a5 move it} & True & 26.4 \\
19 & \texttt{Killing spinors of S\textasciicircum{}4 (radius r0, frame e\textasciicircum{}a = (r0, r0 sin x4, ...) dx\textasciicircum{}(a+3), gt\_\allowbreak{}a = I rho\_\allowbreak{}(a+3)): eps\_\allowbreak{}s = exp(s (i/\allowbreak{}2) x4 gt\_\allowbreak{}1) exp((1/\allowbreak{}2) x5 gt\_\allowbreak{}1 gt\_\allowbreak{}2) exp((1/\allowbreak{}2) x6 gt\_\allowbreak{}2 gt\_\allowbreak{}3) exp((1/\allowbreak{}2) x7 gt\_\allowbreak{}3 gt\_\allowbreak{}4) eps0 satisfies nabla\_\allowbreak{}mu eps\_\allowbreak{}s == s (i/\allowbreak{}(2 r0)) e\_\allowbreak{}mu\textasciicircum{}a gt\_\allowbreak{}a eps\_\allowbreak{}s (s = +-1, all eps0), so the Dirac operator gives (2 s i/\allowbreak{}r0) eps\_\allowbreak{}s: |lambda| = 2/\allowbreak{}r0, Friedrich's bound n R/\allowbreak{}(4(n - 1)) = 4/\allowbreak{}r0\textasciicircum{}2} & True & 26.4 \\
20 & \texttt{Dirac spectrum on S\textasciicircum{}4: Weyl's law for spinors (4 components) N(|lambda| <= L/\allowbreak{}r) \textasciitilde{} 4 V (L/\allowbreak{}r)\textasciicircum{}4/\allowbreak{}(32 Pi\textasciicircum{}2) = L\textasciicircum{}4/\allowbreak{}3 agrees with the degeneracies 2 x 4 Binomial[l + 3, 3] of the eigenvalues +-i(l + 2)/\allowbreak{}r (to 2.0 10\textasciicircum{}-3 at L = 1000)} & True & 26.4 \\
21 & \texttt{the sphere with the volume of the torus today (V = M\_\allowbreak{}Pl\textasciicircum{}2/\allowbreak{}M8\textasciicircum{}6): r = (3 V/\allowbreak{}(8 Pi\textasciicircum{}2))\textasciicircum{}(1/\allowbreak{}4) = 688.9 GeV\textasciicircum{}-1 = 1.359 10\textasciicircum{}-13 m; lowest scalar levels 4.03 MeV (torus, 8 states) and 2.90 MeV (sphere, l = 1, 5 states); lowest dirac16complex level on the sphere 2/\allowbreak{}r = 2.90 MeV (8 Dirac fermions; the torus has 4 massless ones for m = 0); levels below 5 GeV: torus 1.17335 10\textasciicircum{}13, sphere 1.17375 10\textasciicircum{}13, Weyl 1.17335 10\textasciicircum{}13} & True & 26.4 \\
22 & \texttt{RS1 in five dimensions (mostly minus): G\textasciicircum{}mu\_\allowbreak{}nu = 3 A'' - 6 A'\textasciicircum{}2, G\textasciicircum{}y\_\allowbreak{}y = -6 A'\textasciicircum{}2; bulk A = k y needs Lambda = -6 k\textasciicircum{}2; the delta functions A'' = +-2 k delta at the branes need the tensions f\_\allowbreak{}UV = +6 k/\allowbreak{}kappa, f\_\allowbreak{}IR = -6 k/\allowbreak{}kappa (UV positive)} & True & 26.5 \\
23 & \texttt{Israel junction at the UV brane of RS1: with K\textasciicircum{}mu\_\allowbreak{}nu = (1/\allowbreak{}2) g\textasciicircum{}mu alpha d\_\allowbreak{}y g\_\allowbreak{}alpha nu = -A' delta\textasciicircum{}mu\_\allowbreak{}nu (normal d\_\allowbreak{}y), [K\textasciicircum{}mu\_\allowbreak{}nu] - delta\textasciicircum{}mu\_\allowbreak{}nu [K] = 6 k = (coefficient of A'' in G\textasciicircum{}mu\_\allowbreak{}nu) x (jump of A'), the delta-function part of G\textasciicircum{}mu\_\allowbreak{}nu; in the signature (-,+,+,+,+) of the note G\textasciicircum{}A\_\allowbreak{}B changes sign, so there [K\_\allowbreak{}mu nu - g\_\allowbreak{}mu nu K] = -kappa S\_\allowbreak{}mu nu (the note uses +kappa S\_\allowbreak{}mu nu, hence its reversed tensions)} & True & 26.5 \\
24 & \texttt{RS1: M\_\allowbreak{}Pl\textasciicircum{}2 = (1/\allowbreak{}kappa) Int\_\allowbreak{}(-Pi R)\textasciicircum{}(Pi R) e\textasciicircum{}(-2 k |y|) dy = (1 - e\textasciicircum{}(-2 Pi k R))/\allowbreak{}(kappa k)} & True & 26.5 \\
25 & \texttt{RS1 radion: with R promoted to T(x), the bulk action is (3/\allowbreak{}(kappa k)) (d e\textasciicircum{}(-k Pi T))\textasciicircum{}2 + V(T) + (a total derivative), so the canonical radion is mu = Sqrt[6/\allowbreak{}(kappa k)] e\textasciicircum{}(-k Pi T) (the note: Sqrt[12/\allowbreak{}(kappa k)])} & True & 26.5 \\
26 & \texttt{eight dimensions, warped along one extra direction (ds\textasciicircum{}2 = e\textasciicircum{}(-2 A(y)) eta - dy\textasciicircum{}2 - dz.dz): G\textasciicircum{}mu\_\allowbreak{}nu = 3 A'' - 6 A'\textasciicircum{}2, G\textasciicircum{}y\_\allowbreak{}y = -6 A'\textasciicircum{}2, G\textasciicircum{}z\_\allowbreak{}z = 4 A'' - 10 A'\textasciicircum{}2; with only Lambda (all G\textasciicircum{}A\_\allowbreak{}A equal) A'' = 0 and A' = 0: the note's five-dimensional equations do not carry over} & True & 26.5 \\
27 & \texttt{Goldberger-Wise in RS1: the static bulk scalar (Sqrt|G| = e\textasciicircum{}(-4 k y), G\textasciicircum{}yy = -1) obeys Phi'' - 4 k Phi' - m\textasciicircum{}2 Phi = 0, solved by e\textasciicircum{}(2 k y)(C1 e\textasciicircum{}(nu k y) + C2 e\textasciicircum{}(-nu k y)), nu = Sqrt[4 + m\textasciicircum{}2/\allowbreak{}k\textasciicircum{}2] = 2 + eps + O(m\textasciicircum{}4), eps = m\textasciicircum{}2/\allowbreak{}(4 k\textasciicircum{}2)} & True & 26.5 \\
28 & \texttt{Goldberger-Wise radion mass: V(mu) = 4 k (mu/\allowbreak{}f)\textasciicircum{}4 (v\_\allowbreak{}B - v\_\allowbreak{}A (mu/\allowbreak{}f)\textasciicircum{}eps)\textasciicircum{}2 has V'(<mu>) = 0 at (<mu>/\allowbreak{}f)\textasciicircum{}eps = v\_\allowbreak{}B/\allowbreak{}v\_\allowbreak{}A (k Pi <T> = (4 k\textasciicircum{}2/\allowbreak{}m\textasciicircum{}2) Log[v\_\allowbreak{}A/\allowbreak{}v\_\allowbreak{}B]) and V''(<mu>) = 8 k eps\textasciicircum{}2 v\_\allowbreak{}B\textasciicircum{}2 <mu>\textasciicircum{}2/\allowbreak{}f\textasciicircum{}4; with f\textasciicircum{}2 = 6/\allowbreak{}(kappa k): m\_\allowbreak{}r\textasciicircum{}2 = (4/\allowbreak{}3) kappa k\textasciicircum{}2 eps\textasciicircum{}2 v\_\allowbreak{}B\textasciicircum{}2 (<mu>/\allowbreak{}f)\textasciicircum{}2 (the note's f\textasciicircum{}2 = 12/\allowbreak{}(kappa k) gives 2/\allowbreak{}3)} & True & 26.5 \\
29 & \texttt{conformally flat brane metric h = Omega\textasciicircum{}2 eta (four dimensions, Omega(x0, x1)): Sqrt|h| R(h) == -6 Omega box Omega with box = d0\textasciicircum{}2 - d1\textasciicircum{}2 (curvature conventions of 7+1.nb)} & True & 26.6 \\
30 & \texttt{the Higgs boson on the brane (h = Omega\textasciicircum{}2 eta, Omega = 1 + r/\allowbreak{}Lambda + a2 r\textasciicircum{}2, phi\_\allowbreak{}H = v + h, gamma = v/\allowbreak{}Lambda): the terms linear in (h, r) are total derivatives, and to second order L == (1/\allowbreak{}2)(dh)\textasciicircum{}2 + (1/\allowbreak{}2)(1 + 6 xi gamma\textasciicircum{}2)(dr)\textasciicircum{}2 + 6 xi gamma dh.dr - (1/\allowbreak{}2) m\_\allowbreak{}h0\textasciicircum{}2 h\textasciicircum{}2 - (1/\allowbreak{}2) m\_\allowbreak{}r0\textasciicircum{}2 r\textasciicircum{}2 up to total derivatives (Euler-Lagrange expressions in x0 and x1); a2 drops out} & True & 26.6 \\
31 & \texttt{canonical kinetic terms: Transpose[S].K.S == 1 for S = \{\{1, -6 xi gamma/\allowbreak{}Z\}, \{0, 1/\allowbreak{}Z\}\} if and only if Z\textasciicircum{}2 = 1 + 6 xi gamma\textasciicircum{}2 (1 - 6 xi) (the note's Z); with the note's literal kinetic terms (no 6 xi gamma\textasciicircum{}2 for the radion) it would be Z\textasciicircum{}2 = 1 - 36 xi\textasciicircum{}2 gamma\textasciicircum{}2} & True & 26.6 \\
32 & \texttt{mass mixing: with (h', r') = Q(theta).(h\_\allowbreak{}m, r\_\allowbreak{}m), Q = \{\{cos, -sin\}, \{sin, cos\}\}, the off-diagonal element of Transpose[Q].M'.Q is -(1/\allowbreak{}2)(M'\_\allowbreak{}11 - M'\_\allowbreak{}22) sin 2 theta + M'\_\allowbreak{}12 cos 2 theta, zero for tan 2 theta = 2 M'\_\allowbreak{}12/\allowbreak{}(M'\_\allowbreak{}11 - M'\_\allowbreak{}22) = 12 gamma xi Z m\_\allowbreak{}h0\textasciicircum{}2/\allowbreak{}(m\_\allowbreak{}r0\textasciicircum{}2 - m\_\allowbreak{}h0\textasciicircum{}2 (Z\textasciicircum{}2 - 36 xi\textasciicircum{}2 gamma\textasciicircum{}2)) (the note's formula); with the inverse matrix (the note's display) it is +(1/\allowbreak{}2)(M'\_\allowbreak{}11 - M'\_\allowbreak{}22) sin 2 theta + M'\_\allowbreak{}12 cos 2 theta (theta -> -theta); (h, r) = S.Q.(h\_\allowbreak{}m, r\_\allowbreak{}m) gives C\_\allowbreak{}h = cos theta - (6 xi gamma/\allowbreak{}Z) sin theta, C\_\allowbreak{}r = sin theta/\allowbreak{}Z} & True & 26.6 \\
33 & \texttt{Higgs-radion mixing in 7+1 (gamma = v/\allowbreak{}(Sqrt[3] M\_\allowbreak{}Pl), xi = 1, m\_\allowbreak{}r0 = 10 GeV, m\_\allowbreak{}h0 = 125.20 GeV; 60-digit arithmetic): |Z - 1| < 10\textasciicircum{}-30, |theta| = 3.5 10\textasciicircum{}-16, the rates of h -> gamma gamma and h -> g g change by less than 10\textasciicircum{}-30; in RS1 with <mu> = 3 TeV gamma = v/\allowbreak{}(3 TeV) = 0.082} & True & 26.6 \\*
\end{longtable}}
\Needspace{8\baselineskip}\subsection*{6.2 The independent Python programs}\addcontentsline{toc}{subsection}{6.2 The independent Python programs}
tools/ref\_torus\_brane.py — its boolean checks:

{\footnotesize
\begin{longtable}{>{\raggedright\arraybackslash}p{\dimexpr 0.0550\linewidth-2\tabcolsep\relax}>{\raggedright\arraybackslash}p{\dimexpr 0.8250\linewidth-2\tabcolsep\relax}>{\raggedright\arraybackslash}p{\dimexpr 0.1200\linewidth-2\tabcolsep\relax}}
\toprule
\textbf{\#} & \textbf{check} & \textbf{result} \\
\midrule\endfirsthead
\toprule
\textbf{\#} & \textbf{check} & \textbf{result} \\
\midrule\endhead
\endfoot
\bottomrule\endlastfoot
1 & \texttt{torus: det(-g\_\allowbreak{}kl) = b\textasciicircum{}8 for the cosmological metric; the proper volume of [0, Lc)\textasciicircum{}4 is b\textasciicircum{}4 Lc\textasciicircum{}4} & True \\
2 & \texttt{only b Lc is physical: y = lambda y' maps (Lc, b) to (Lc/\allowbreak{}lambda, lambda b); b\textasciicircum{}2 dy.dy, b Lc and b\textasciicircum{}4 Lc\textasciicircum{}4 unchanged} & True \\
3 & \texttt{symmetries fixing y = 0: 384 = 2\textasciicircum{}4 4! distinct signed permutation matrices, orthogonal integer matrices; shells 1, 2, 3, 4, 9 have orbits \{8\}, \{24\}, \{32\}, \{8, 16\}, \{8, 96\}} & True \\
4 & \texttt{Fourier modes phi\_\allowbreak{}n = exp(i q\_\allowbreak{}n.y)/\allowbreak{}sqrt(V\_\allowbreak{}perp), q\_\allowbreak{}n = 2 pi n/\allowbreak{}Lc: orthonormal for b\textasciicircum{}4 d\textasciicircum{}4y (three exact pairs)} & True \\
5 & \texttt{Parseval: f = prod cos\textasciicircum{}2(pi y\_\allowbreak{}k/\allowbreak{}Lc) has 81 coefficients sqrt(V\_\allowbreak{}perp) prod w(n\_\allowbreak{}k) (3 checked by integration), Sum |c\_\allowbreak{}n|\textasciicircum{}2 = Int b\textasciicircum{}4 f\textasciicircum{}2 = (3/\allowbreak{}8)\textasciicircum{}4 V\_\allowbreak{}perp} & True \\
6 & \texttt{Laplacian of the torus: -(1/\allowbreak{}b\textasciicircum{}2) Sum d\textasciicircum{}2/\allowbreak{}dy\_\allowbreak{}k\textasciicircum{}2 phi\_\allowbreak{}n = (|q\_\allowbreak{}n|\textasciicircum{}2/\allowbreak{}b\textasciicircum{}2) phi\_\allowbreak{}n} & True \\
7 & \texttt{scalar Kaluza-Klein masses: (Box + m\textasciicircum{}2) exp(-iEx0 + ik.x + iq.y) = (-E\textasciicircum{}2 + |k|\textasciicircum{}2 + |q|\textasciicircum{}2/\allowbreak{}b0\textasciicircum{}2 + m\textasciicircum{}2) exp(...) on diag(1,-1,-1,-1,-b0\textasciicircum{}2,...)} & True \\
8 & \texttt{graviton on the torus: D(D-3)/\allowbreak{}2 = 20 (D = 8) = 2 + 8 + 10 (zero mode: graviton, 4 graviphotons, 10 moduli) = 5 + 9 + 6 (massive level)} & True \\
9 & \texttt{Jacobi: \#\{n in Z\textasciicircum{}4 : |n|\textasciicircum{}2 = k\} = r\_\allowbreak{}4(k) = 8 sigma(k) - 32 sigma(k/\allowbreak{}4), direct count for k = 1..100} & True \\
10 & \texttt{Weyl's law: N(10) = 49689, N(100) = 493490641; |N(R)/\allowbreak{}(pi\textasciicircum{}2 R\textasciicircum{}4/\allowbreak{}2) - 1| < 1e-2, 1e-3, 3e-5, 5e-6, 1e-6} & True \\
11 & \texttt{numbers today: V\_\allowbreak{}perp = 5.929e12 GeV\textasciicircum{}-4, l\_\allowbreak{}perp = 1560.4 GeV\textasciicircum{}-1 = 3.079e-13 m, 2 pi/\allowbreak{}l\_\allowbreak{}perp = 4.027 MeV; N(5 GeV) = 1.17e13 (lattice to 1e-5); N(M8) = M\_\allowbreak{}Pl\textasciicircum{}2/\allowbreak{}(32 pi\textasciicircum{}2 M8\textasciicircum{}2)} & True \\
12 & \texttt{from eight to four dimensions: for diag(1, -a\textasciicircum{}2, -a\textasciicircum{}2, -a\textasciicircum{}2, -b0\textasciicircum{}2, ..., -b0\textasciicircum{}2) (static torus) R = R4 = -6(Ha' + 2 Ha\textasciicircum{}2), the scalar curvature of the FRW metric diag(1, -a\textasciicircum{}2, -a\textasciicircum{}2, -a\textasciicircum{}2); sqrt|g| = a\textasciicircum{}3 b0\textasciicircum{}4} & True \\
13 & \texttt{quaternion subalgebra of the brane directions: e1 e2 = -e3, e2 e3 = -e1, e3 e1 = -e2 (e\_\allowbreak{}i e\_\allowbreak{}j = -delta\_\allowbreak{}ij - eps\_\allowbreak{}ijk e\_\allowbreak{}k), span\{e0, ..., e3\} closed under the octonion product, Re(e\_\allowbreak{}A e\_\allowbreak{}B) = diag(1, -1, -1, -1)} & True \\
14 & \texttt{Newton in 8D: Omega\_\allowbreak{}6 = 16 pi\textasciicircum{}3/\allowbreak{}15, (1/\allowbreak{}r\textasciicircum{}6) d/\allowbreak{}dr (r\textasciicircum{}6 d/\allowbreak{}dr r\textasciicircum{}-5) = 0, Phi = -(5/\allowbreak{}6) kappa M/\allowbreak{}(5 Omega\_\allowbreak{}6 r\textasciicircum{}5) = -5 kappa M/\allowbreak{}(32 pi\textasciicircum{}3 r\textasciicircum{}5); D = 4: -G M/\allowbreak{}r} & True \\
15 & \texttt{4D Fourier transform: Int d\textasciicircum{}4y exp(-ip.y) (r\textasciicircum{}2 + |y|\textasciicircum{}2)\textasciicircum{}(-5/\allowbreak{}2) = (4 pi\textasciicircum{}2/\allowbreak{}(3 r)) exp(-p r) (mpmath, three (p, r), to 1e-20)} & True \\
16 & \texttt{Poisson duality: Sum\_\allowbreak{}n (r\textasciicircum{}2 + |n|\textasciicircum{}2)\textasciicircum{}(-5/\allowbreak{}2) = (4 pi\textasciicircum{}2/\allowbreak{}(3 r)) Sum\_\allowbreak{}m exp(-2 pi r |m|) to 1e-8 at r = 1/\allowbreak{}4, 1/\allowbreak{}2, 1, 2} & True \\
17 & \texttt{Z(5/\allowbreak{}2) = 7 zeta(5/\allowbreak{}2) zeta(3/\allowbreak{}2) = Sum r\_\allowbreak{}4(k) k\textasciicircum{}-5/\allowbreak{}2 (to 1e-8); near field ratio -> 1 (0.98 < . < 1 at r = 1/\allowbreak{}10, 1/\allowbreak{}20); far field ratio -> 1 (to 2e-2 at r = 2, 3)} & True \\
18 & \texttt{graviton + radion: (5 kappa M/\allowbreak{}(32 pi\textasciicircum{}3))(4 pi\textasciicircum{}2/\allowbreak{}(3 V\_\allowbreak{}perp r)) = (5/\allowbreak{}3) G4 M/\allowbreak{}r; radion (1/\allowbreak{}(sqrt3 M\_\allowbreak{}Pl))\textasciicircum{}2/\allowbreak{}(4 pi) = (2/\allowbreak{}3)/\allowbreak{}(8 pi M\_\allowbreak{}Pl\textasciicircum{}2)} & True \\
19 & \texttt{Kaluza-Klein momentum: bulk vertex Int b\textasciicircum{}4 phi\_\allowbreak{}n phi\_\allowbreak{}m phi\_\allowbreak{}l\textasciicircum{}* = 1/\allowbreak{}sqrt(V\_\allowbreak{}perp) for l = n + m, 0 otherwise; at y = 0 phi\_\allowbreak{}n phi\_\allowbreak{}m phi\_\allowbreak{}l\textasciicircum{}* = V\_\allowbreak{}perp\textasciicircum{}(-3/\allowbreak{}2)} & True \\
20 & \texttt{delta function of the torus: Fourier coefficients Int exp(-i q\_\allowbreak{}n y) delta(y) dy = 1; Int b\textasciicircum{}4 d\textasciicircum{}4y rho4 delta\textasciicircum{}4(y)/\allowbreak{}b\textasciicircum{}4 = rho4 (average rho4/\allowbreak{}V\_\allowbreak{}perp)} & True \\
21 & \texttt{induced metric h = J\textasciicircum{}T g J: FRW diag(1, -a\textasciicircum{}2, -a\textasciicircum{}2, -a\textasciicircum{}2) for the cosmological metric; diag(n\textasciicircum{}2, -a\textasciicircum{}2, -a\textasciicircum{}2, -a\textasciicircum{}2) at x4 = 0 with n = 1 + x0\textasciicircum{}2/\allowbreak{}11, a = 1 + x0/\allowbreak{}3 for the inhomogeneous metric} & True \\
22 & \texttt{brane volume element: sqrt|g| = sqrt|h| b\textasciicircum{}4 at the brane (cosmological metric; inhomogeneous metric at x0 = 1/\allowbreak{}2, exact), so Int\_\allowbreak{}W d\textasciicircum{}4x sqrt|h| L = Int d\textasciicircum{}8X sqrt|g| delta\textasciicircum{}4(y)/\allowbreak{}b\textasciicircum{}4 L; test Int b\textasciicircum{}4 (delta\textasciicircum{}4/\allowbreak{}b\textasciicircum{}4) F = F(0)} & True \\
23 & \texttt{Christoffel symbols \{k mu nu\} at the brane: all 0 for the cosmological metric (totally geodesic); for the inhomogeneous metric at x0 = 1/\allowbreak{}2 only \{4 00\} = 1458/\allowbreak{}58619 = n d4n/\allowbreak{}b\textasciicircum{}2} & True \\
24 & \texttt{R\_\allowbreak{}A = (gamma\_\allowbreak{}mu (x) 1\_\allowbreak{}4, gamma5 (x) rho\_\allowbreak{}k) and Gamma17[A] both satisfy X\_\allowbreak{}A X\_\allowbreak{}B + X\_\allowbreak{}B X\_\allowbreak{}A = 2 eta\_\allowbreak{}AB 1\_\allowbreak{}16} & True \\
25 & \texttt{every Gamma17[A] is a signed permutation matrix} & True \\
26 & \texttt{intertwiner: Schur average with the seed matrix unit (1, 2); 2 sqrt(2) U has entries 0, +-1, +-i (rounding error < 1e-12) with 8 non-zero entries in each row and column; Gamma17[A].U8 = U8.R\_\allowbreak{}A and U8\textasciicircum{}dag.U8 = 8 1\_\allowbreak{}16 exactly} & True \\
27 & \texttt{Schur's lemma: the linear map X -> (Gamma\_\allowbreak{}A X - X R\_\allowbreak{}A)\_\allowbreak{}A has a one-dimensional null space (one singular value < 1e-10, the next > 0.1)} & True \\
28 & \texttt{brane Lorentz group on the brane spinor index, transverse rotations on a: Gamma\_\allowbreak{}mu Gamma\_\allowbreak{}nu = U (gamma\_\allowbreak{}mu gamma\_\allowbreak{}nu (x) 1) U\textasciicircum{}dag, Gamma\_\allowbreak{}k Gamma\_\allowbreak{}l = U (1 (x) rho\_\allowbreak{}k rho\_\allowbreak{}l) U\textasciicircum{}dag, Gamma\_\allowbreak{}mu Gamma\_\allowbreak{}k = U (gamma\_\allowbreak{}mu gamma5 (x) rho\_\allowbreak{}k) U\textasciicircum{}dag} & True \\
29 & \texttt{Spin(4) on the transverse index: rho4567 = -gamma5 (4x4) commutes with rho\_\allowbreak{}k rho\_\allowbreak{}l, anticommutes with rho\_\allowbreak{}k; (rho\_\allowbreak{}k rho\_\allowbreak{}l + s rho\_\allowbreak{}m rho\_\allowbreak{}n) P\_\allowbreak{}s = 0 and (rho\_\allowbreak{}k rho\_\allowbreak{}l - s rho\_\allowbreak{}m rho\_\allowbreak{}n) P\_\allowbreak{}s has rank 2 for even (k,l,m,n): 4 = (2,1) + (1,2)} & True \\
30 & \texttt{chirality: chi = i Gamma\_\allowbreak{}0...Gamma\_\allowbreak{}7 = U (gamma5 (x) rho4567) U\textasciicircum{}dag; P\_\allowbreak{}L17 = U (P\_\allowbreak{}L (x) P\_\allowbreak{}+ + P\_\allowbreak{}R (x) P\_\allowbreak{}-) U\textasciicircum{}dag, each part of rank 4} & True \\
31 & \texttt{example psi = U.(u (x) e\_\allowbreak{}1), u = (1, 2, 0, -1): psibar psi = ubar u = 4 and psibar Gamma\textasciicircum{}mu psi = ubar gamma\textasciicircum{}mu u} & True \\
32 & \texttt{four brane Dirac spinors: psibar psi' = Sum\_\allowbreak{}a psibar\_\allowbreak{}a psi'\_\allowbreak{}a and psibar Gamma\textasciicircum{}mu psi' = Sum\_\allowbreak{}a psibar\_\allowbreak{}a gamma\textasciicircum{}mu psi'\_\allowbreak{}a (random integer spinors)} & True \\
33 & \texttt{Kaluza-Klein operator X = gamma5 (x) (rho.q)/\allowbreak{}b0: X\textasciicircum{}2 = -|q|\textasciicircum{}2/\allowbreak{}b0\textasciicircum{}2 1\_\allowbreak{}16 and X anticommutes with gamma\_\allowbreak{}mu (x) 1\_\allowbreak{}4} & True \\
34 & \texttt{brane tension: T\_\allowbreak{}mu nu = (2/\allowbreak{}sqrt|h|) dS/\allowbreak{}dh\textasciicircum{}mu nu = f\_\allowbreak{}W h\_\allowbreak{}mu nu for S = -Int sqrt|h| f\_\allowbreak{}W; T\textasciicircum{}0\_\allowbreak{}0 = f\_\allowbreak{}W, T\textasciicircum{}1\_\allowbreak{}1 = f\_\allowbreak{}W (p = -f\_\allowbreak{}W, w = -1)} & True \\
35 & \texttt{smeared sources: tension eps = f\_\allowbreak{}W/\allowbreak{}V\_\allowbreak{}perp, pa = -eps, pb = 0; brane radiation pa = eps/\allowbreak{}3, pb = 0; bulk radiation <cos\textasciicircum{}2> = 3/\allowbreak{}7, <sin\textasciicircum{}2> = 4/\allowbreak{}7: pa = pb = eps/\allowbreak{}7} & True \\
36 & \texttt{normalcy temperature: (T\textasciicircum{}7/\allowbreak{}M8\textasciicircum{}6)/\allowbreak{}(T\textasciicircum{}2/\allowbreak{}M\_\allowbreak{}Pl) = 1 at T* = (M8\textasciicircum{}6/\allowbreak{}M\_\allowbreak{}Pl)\textasciicircum{}(1/\allowbreak{}5) = 13.3 GeV; V\_\allowbreak{}perp T\textasciicircum{}4 T\textasciicircum{}3/\allowbreak{}M\_\allowbreak{}Pl\textasciicircum{}2 = T\textasciicircum{}7/\allowbreak{}M8\textasciicircum{}6} & True \\
37 & \texttt{the tension is vacuum energy: kappa f\_\allowbreak{}W/\allowbreak{}V\_\allowbreak{}perp = f\_\allowbreak{}W/\allowbreak{}M\_\allowbreak{}Pl\textasciicircum{}2 (V\_\allowbreak{}perp = kappa M\_\allowbreak{}Pl\textasciicircum{}2); f\_\allowbreak{}W = M8\textasciicircum{}4 is 3e62 to 5e62 rho\_\allowbreak{}Lambda} & True \\
38 & \texttt{Green's function of the four transverse directions: (1/\allowbreak{}y\textasciicircum{}3) d/\allowbreak{}dy (y\textasciicircum{}3 d/\allowbreak{}dy y\textasciicircum{}-2) = 0, flux 2 pi\textasciicircum{}2 y\textasciicircum{}3 d/\allowbreak{}dy y\textasciicircum{}-2 = -4 pi\textasciicircum{}2; sqrt(kappa f\_\allowbreak{}W) = 1/\allowbreak{}M8 for f\_\allowbreak{}W = M8\textasciicircum{}4} & True \\*
\end{longtable}}
tools/ref\_radion.py — its boolean checks:

{\footnotesize
\begin{longtable}{>{\raggedright\arraybackslash}p{\dimexpr 0.0550\linewidth-2\tabcolsep\relax}>{\raggedright\arraybackslash}p{\dimexpr 0.8250\linewidth-2\tabcolsep\relax}>{\raggedright\arraybackslash}p{\dimexpr 0.1200\linewidth-2\tabcolsep\relax}}
\toprule
\textbf{\#} & \textbf{check} & \textbf{result} \\
\midrule\endfirsthead
\toprule
\textbf{\#} & \textbf{check} & \textbf{result} \\
\midrule\endhead
\endfoot
\bottomrule\endlastfoot
1 & \texttt{the 4-sphere of radius r0: Ricci\_\allowbreak{}ij = 3 gamma\_\allowbreak{}ij, off-diagonal 0, R = -12/\allowbreak{}r0\textasciicircum{}2 in the signature of 7+1.nb; volume of the unit S\textasciicircum{}4 = 8 pi\textasciicircum{}2/\allowbreak{}3} & True \\
2 & \texttt{radion on T\textasciicircum{}4: the reduced action equals -(M\_\allowbreak{}Pl\textasciicircum{}2/\allowbreak{}2) R\_\allowbreak{}E + 6 M\_\allowbreak{}Pl\textasciicircum{}2 (d beta)\textasciicircum{}2 - e\textasciicircum{}(-4 beta)(M\_\allowbreak{}Pl\textasciicircum{}2 Lambda + V0 U) up to a total derivative (Euler-Lagrange for N, a, beta), M\_\allowbreak{}Pl\textasciicircum{}2 = V0/\allowbreak{}kappa} & True \\
3 & \texttt{radion on S\textasciicircum{}4: the same with the potential -(M\_\allowbreak{}Pl\textasciicircum{}2/\allowbreak{}2) e\textasciicircum{}(-4 beta) 12/\allowbreak{}(r0\textasciicircum{}2 e\textasciicircum{}(2 beta)) from the curvature of the sphere, V0 = (8 pi\textasciicircum{}2/\allowbreak{}3) r0\textasciicircum{}4} & True \\
4 & \texttt{n extra dimensions (T\textasciicircum{}n, n = 1..4): kinetic term (M\_\allowbreak{}Pl\textasciicircum{}2/\allowbreak{}2)(n(n+2)/\allowbreak{}2)(dB)\textasciicircum{}2; the note's beta\_\allowbreak{}n = sqrt(2/\allowbreak{}(n(n+2))) gives (M\_\allowbreak{}Pl\textasciicircum{}2/\allowbreak{}2)(d phi)\textasciicircum{}2, not (M\_\allowbreak{}Pl\textasciicircum{}2/\allowbreak{}2)(1/\allowbreak{}2)(d phi)\textasciicircum{}2; beta\_\allowbreak{}n = 1/\allowbreak{}sqrt(n(n+2)) gives the 1/\allowbreak{}2} & True \\
5 & \texttt{the note's eight-dimensional alpha = -1/\allowbreak{}sqrt(3), beta = 1/\allowbreak{}sqrt(12): 2 alpha + 4 beta = 0, 4 beta = 2/\allowbreak{}sqrt(3), sqrt(12) beta = 1} & True \\
6 & \texttt{the kinetic coefficient solved from the reduction (n = 1..4) is n(n+2)/\allowbreak{}2} & True \\
7 & \texttt{the cosmological metric with the sphere of radius b: G\textasciicircum{}A\_\allowbreak{}A = torus + (6, 6, 6, 6, 3, 3, 3, 3)/\allowbreak{}b\textasciicircum{}2, R = torus - 12/\allowbreak{}b\textasciicircum{}2} & True \\
8 & \texttt{radion equation on S\textasciicircum{}4: Hb' + 3 Ha Hb + 4 Hb\textasciicircum{}2 = kappa (eps/\allowbreak{}6 - pa/\allowbreak{}2 + pb/\allowbreak{}3) + Lambda/\allowbreak{}3 - 3/\allowbreak{}b\textasciicircum{}2} & True \\
9 & \texttt{radion on T\textasciicircum{}4 with U: flat vacuum Lambda + kappa U0 = 0, U1 = 0; m\_\allowbreak{}r\textasciicircum{}2 = kappa K\_\allowbreak{}U/\allowbreak{}12} & True \\
10 & \texttt{radion on S\textasciicircum{}4 with U: Lambda + kappa U0 = 6/\allowbreak{}r0\textasciicircum{}2, kappa U1 = -12/\allowbreak{}r0\textasciicircum{}2, m\_\allowbreak{}r\textasciicircum{}2 = kappa K\_\allowbreak{}U/\allowbreak{}12 - 2/\allowbreak{}r0\textasciicircum{}2, from V and from the 8D radion equation} & True \\
11 & \texttt{S\textasciicircum{}4 with Lambda only: dS4 x S\textasciicircum{}4 (r0\textasciicircum{}2 = 9/\allowbreak{}Lambda, H\textasciicircum{}2 = 1/\allowbreak{}r0\textasciicircum{}2) solves the 00, 11, 44 equations; maximum of V with m\textasciicircum{}2 = -6 H\textasciicircum{}2, also from the 8D radion equation} & True \\
12 & \texttt{S\textasciicircum{}4 with a 4-form flux: eps = rho\_\allowbreak{}F, pa = -rho\_\allowbreak{}F, pb = +rho\_\allowbreak{}F; Minkowski needs Lambda = 9/\allowbreak{}(2 r0\textasciicircum{}2), kappa rho\_\allowbreak{}F = 3/\allowbreak{}(2 r0\textasciicircum{}2); m\textasciicircum{}2 = 6/\allowbreak{}r0\textasciicircum{}2 from V and from the 8D radion equation} & True \\
13 & \texttt{the note's Goldberger-Wise potential: for a static Phi(y) the Einstein-frame energy carries e\textasciicircum{}(-4 beta) (gradient part e\textasciicircum{}(-6 beta)), not e\textasciicircum{}(+4 beta)} & True \\
14 & \texttt{radion coupling to a mass on the brane: -M e\textasciicircum{}(-n beta/\allowbreak{}2) = -M + (n/\allowbreak{}2) M beta; coupling sqrt(n/\allowbreak{}(2(n+2)))/\allowbreak{}M\_\allowbreak{}Pl = 1/\allowbreak{}(sqrt(3) M\_\allowbreak{}Pl) for n = 4; fifth force 2 (coupling M\_\allowbreak{}Pl)\textasciicircum{}2 = n/\allowbreak{}(n+2) = 2/\allowbreak{}3} & True \\
15 & \texttt{Eot-Wash: m\_\allowbreak{}r > hbar c/\allowbreak{}38.6 micrometres = 5.1 meV; range of a 10 GeV radion 2.0e-17 m; gamma = v/\allowbreak{}(sqrt(3) M\_\allowbreak{}Pl) = 5.8e-17} & True \\
16 & \texttt{zonal harmonics on S\textasciicircum{}4 (l = 0..6): eigenvalue -l(l+3), C\_\allowbreak{}l(1) = (l+1)(l+2)/\allowbreak{}2, Int (C\_\allowbreak{}l/\allowbreak{}C\_\allowbreak{}l(1))\textasciicircum{}2 = V/\allowbreak{}d\_\allowbreak{}l; d\_\allowbreak{}l = binomial(l+4,4) - binomial(l+2,4) (l <= 20)} & True \\
17 & \texttt{Weyl's law on S\textasciicircum{}4: N(x)/\allowbreak{}(x\textasciicircum{}4/\allowbreak{}12) - 1 = -1/\allowbreak{}x\textasciicircum{}2 at x = 10, 100, 1000} & True \\
18 & \texttt{the brane at the north pole of S\textasciicircum{}4: the embedding induces r0\textasciicircum{}2 gamma; north pole (0,0,0,0,r0); 6 of the 10 rotations of R\textasciicircum{}5 fix it (SO(4)), 4 move it} & True \\
19 & \texttt{Killing spinors of S\textasciicircum{}4: nabla eps\_\allowbreak{}s = s (i/\allowbreak{}(2 r0)) gt eps\_\allowbreak{}s and D eps\_\allowbreak{}s = (2 s i/\allowbreak{}r0) eps\_\allowbreak{}s (s = +-1; residual < 1e-25 at four random points, 30 digits); Friedrich bound 4/\allowbreak{}r0\textasciicircum{}2} & True \\
20 & \texttt{Dirac spectrum on S\textasciicircum{}4: Weyl's law for spinors L\textasciicircum{}4/\allowbreak{}3 versus the degeneracies 2 x 4 binomial(l+3,3) (2.0e-3 at L = 1000)} & True \\
21 & \texttt{the sphere with the volume of the torus today: r = 688.9 GeV\textasciicircum{}-1; lowest levels 4.03 MeV (torus) and 2.90 MeV (sphere); dirac16complex on the sphere 2.90 MeV; levels below 5 GeV agree with Weyl} & True \\
22 & \texttt{RS1 (5D, mostly minus): G\textasciicircum{}mu\_\allowbreak{}nu = 3 A'' - 6 A'\textasciicircum{}2, G\textasciicircum{}y\_\allowbreak{}y = -6 A'\textasciicircum{}2; bulk Lambda = -6 k\textasciicircum{}2; tensions f\_\allowbreak{}UV = +6k/\allowbreak{}kappa, f\_\allowbreak{}IR = -6k/\allowbreak{}kappa} & True \\
23 & \texttt{Israel junction at the UV brane: K\textasciicircum{}mu\_\allowbreak{}nu = -A' delta, [K\textasciicircum{}mu\_\allowbreak{}nu] - delta [K] = 6k = 3 x (jump of A'), the delta part of G\textasciicircum{}mu\_\allowbreak{}nu; in the signature (-,+,+,+,+) [K - g K] = -kappa S} & True \\
24 & \texttt{RS1: M\_\allowbreak{}Pl\textasciicircum{}2 = (1/\allowbreak{}kappa) Int e\textasciicircum{}(-2k|y|) dy over [-pi R, pi R] = (1 - e\textasciicircum{}(-2 pi k R))/\allowbreak{}(kappa k)} & True \\
25 & \texttt{RS1 radion: kinetic term (3/\allowbreak{}(kappa k)) (d e\textasciicircum{}(-k pi T))\textasciicircum{}2, canonical mu = sqrt(6/\allowbreak{}(kappa k)) e\textasciicircum{}(-k pi T)} & True \\
26 & \texttt{8D warped along one extra direction: G\textasciicircum{}y\_\allowbreak{}y = -6 A'\textasciicircum{}2, G\textasciicircum{}z\_\allowbreak{}z = 4 A'' - 10 A'\textasciicircum{}2; with only Lambda A'' = 0 and A' = 0} & True \\
27 & \texttt{Goldberger-Wise: Phi'' - 4k Phi' - m\textasciicircum{}2 Phi = 0 from the bulk action, solved by e\textasciicircum{}(2ky)(C1 e\textasciicircum{}(nu k y) + C2 e\textasciicircum{}(-nu k y)), nu = sqrt(4 + m\textasciicircum{}2/\allowbreak{}k\textasciicircum{}2) = 2 + eps; V'(<mu>) = 0 at (<mu>/\allowbreak{}f)\textasciicircum{}eps = v\_\allowbreak{}B/\allowbreak{}v\_\allowbreak{}A, V''(<mu>) = 8 k eps\textasciicircum{}2 v\_\allowbreak{}B\textasciicircum{}2 <mu>\textasciicircum{}2/\allowbreak{}f\textasciicircum{}4; m\_\allowbreak{}r\textasciicircum{}2 = (4/\allowbreak{}3) kappa k\textasciicircum{}2 eps\textasciicircum{}2 v\_\allowbreak{}B\textasciicircum{}2 (<mu>/\allowbreak{}f)\textasciicircum{}2 (note: 2/\allowbreak{}3)} & True \\
28 & \texttt{conformally flat h = Omega\textasciicircum{}2 eta: sqrt|h| R(h) = -6 Omega box Omega, box = d0\textasciicircum{}2 - d1\textasciicircum{}2} & True \\
29 & \texttt{the Higgs boson on the brane: linear terms are total derivatives; to second order L = (1/\allowbreak{}2)(dh)\textasciicircum{}2 + (1/\allowbreak{}2)(1 + 6 xi gamma\textasciicircum{}2)(dr)\textasciicircum{}2 + 6 xi gamma dh.dr - (1/\allowbreak{}2) m\_\allowbreak{}h0\textasciicircum{}2 h\textasciicircum{}2 - (1/\allowbreak{}2) m\_\allowbreak{}r0\textasciicircum{}2 r\textasciicircum{}2 up to total derivatives; o\_\allowbreak{}2 drops out} & True \\
30 & \texttt{the shift h = h' - (6 xi gamma/\allowbreak{}Z) r', r = r'/\allowbreak{}Z is canonical iff Z\textasciicircum{}2 = 1 + 6 xi gamma\textasciicircum{}2 (1 - 6 xi) (note's literal kinetic terms: 1 - 36 xi\textasciicircum{}2 gamma\textasciicircum{}2); tan 2 theta = 12 gamma xi Z m\_\allowbreak{}h0\textasciicircum{}2/\allowbreak{}(m\_\allowbreak{}r0\textasciicircum{}2 - m\_\allowbreak{}h0\textasciicircum{}2 (Z\textasciicircum{}2 - 36 xi\textasciicircum{}2 gamma\textasciicircum{}2)) diagonalizes with (h', r') = Q(theta)(h\_\allowbreak{}m, r\_\allowbreak{}m); the inverse matrix needs -theta; C\_\allowbreak{}h = cos - (6 xi gamma/\allowbreak{}Z) sin, C\_\allowbreak{}r = sin/\allowbreak{}Z} & True \\
31 & \texttt{Higgs-radion mixing in 7+1 (gamma = v/\allowbreak{}(sqrt(3) M\_\allowbreak{}Pl), xi = 1, m\_\allowbreak{}r0 = 10 GeV): |Z - 1| < 1e-30, |theta| = 3.5e-16, rates change by < 1e-30; warped gamma = v/\allowbreak{}(3 TeV) = 0.082} & True \\*
\end{longtable}}
\textbf{Numbers printed by the notebook and by the programs} (largest relative difference per group, and the tolerance used by tools/verify\_7p1\_nb.py; integers and exact matrices must agree exactly):

{\footnotesize
\begin{longtable}{>{\raggedright\arraybackslash}p{\dimexpr 0.5200\linewidth-2\tabcolsep\relax}>{\raggedright\arraybackslash}p{\dimexpr 0.1000\linewidth-2\tabcolsep\relax}>{\raggedright\arraybackslash}p{\dimexpr 0.2200\linewidth-2\tabcolsep\relax}>{\raggedright\arraybackslash}p{\dimexpr 0.1600\linewidth-2\tabcolsep\relax}}
\toprule
\textbf{quantity} & \textbf{values} & \textbf{max. relative difference} & \textbf{tolerance} \\
\midrule\endfirsthead
\toprule
\textbf{quantity} & \textbf{values} & \textbf{max. relative difference} & \textbf{tolerance} \\
\midrule\endhead
\endfoot
\bottomrule\endlastfoot
orbit sizes of the shells $k=1,\dots,12$ & 16 & 0 & exact \\
$r_4(k)$, $k=0,\dots,16$ & 17 & 0 & exact \\
Weyl: $N(R)$, $R=10,30,100,300,1000$ & 5 & 0 & exact \\
Weyl: $N(R)/(\pi^2R^4/2)-1$ & 5 & $3.8\times10^{-10}$ & $1\times10^{-9}$ \\
$V_\perp$, $\ell_\perp$ (GeV$^{-1}$, m, fm), $2\pi/\ell_\perp$ & 5 & $2.2\times10^{-16}$ & $1\times10^{-12}$ \\
$\ell_\perp/r_p$, $a_0/\ell_\perp$, $R=\ell_\perp/2\pi$, $M_D$ & 5 & $1.7\times10^{-16}$ & $1\times10^{-12}$ \\
$N(\mu)$ at 5 GeV, $m_h$, $M_8$; lattice count below 5 GeV & 4 & $1.8\times10^{-16}$ & $1\times10^{-12}$ \\
the first levels ($k$, $r_4$, $4r_4$, mass) & 32 & $2.5\times10^{-16}$ & $1\times10^{-12}$ \\
image sums $S(r)$ and Kaluza–Klein sums, $r=\frac14,\frac12,1,2$ & 8 & $1.2\times10^{-15}$ & $1\times10^{-9}$ \\
$Z(5/2)$: closed form, shell sum & 2 & $2.9\times10^{-16}$ & $1\times10^{-9}$ \\
near-field and far-field ratios & 5 & $7.8\times10^{-12}$ & $1\times10^{-6}$ \\
Christoffel symbol $\{4\,00\}$ at the brane & 1 & 0 & $1\times10^{-15}$ \\
$\Gamma_A$ as signed permutations & 128 & 0 & exact \\
$\mathrm{Re}(2\sqrt2\,U)$ & 256 & 0 & exact \\
$\mathrm{Im}(2\sqrt2\,U)$ & 256 & 0 & exact \\
example spinor: $\bar\psi\psi$, $\bar\psi\Gamma^\mu\psi$ & 5 & 0 & $1\times10^{-15}$ \\
$T_*$, $f_W/\rho_\Lambda$, $1/(2\pi M_8)$ & 3 & $5.8\times10^{-16}$ & $1\times10^{-12}$ \\
kinetic coefficient $n(n+2)/2$, $n=1,\dots,4$ (solved from the reduction) & 4 & 0 & exact \\
$\beta_n^2$ of the note and the canonical $\beta_n^2$, $n=4$ & 2 & 0 & $1\times10^{-15}$ \\
$S^4$ with $U$: $(\Lambda+\kappa U_0)r_0^2$, $\kappa U_1r_0^2$, $(m_r^2-\kappa K_U/12)r_0^2$ & 3 & 0 & exact \\
$S^4$ with $\Lambda$ only: $\Lambda r_0^2$, $H^2r_0^2$, $m^2/H^2$ & 3 & 0 & exact \\
$S^4$ with a 4-form flux: $\Lambda r_0^2$, $\kappa\rho_Fr_0^2$, $m^2r_0^2$ & 3 & 0 & $1\times10^{-15}$ \\
E\textbackslash{}"ot-Wash bound, range of a 10 GeV radion, $\gamma=v/(\sqrt3M_{\rm Pl})$ & 3 & $2.1\times10^{-16}$ & $1\times10^{-10}$ \\
multiplicities $d_l$ on $S^4$, $l=0,\dots,8$ & 9 & 0 & exact \\
Weyl's law on $S^4$: $N/(x^4/12)-1$ at $x=10,100,1000$ & 3 & $8.2\times10^{-11}$ & $1\times10^{-9}$ \\
Dirac degeneracies on $S^4$ and the spinor Weyl ratio & 7 & $7.0\times10^{-14}$ & $1\times10^{-9}$ \\
today: $r_{S^4}$, lowest masses, numbers of levels below 5 GeV & 8 & $2.2\times10^{-16}$ & $1\times10^{-12}$ \\
RS1: $3$, $f_{\rm UV}\kappa/k$, $f_{\rm IR}\kappa/k$, $f^2\kappa k$ & 4 & 0 & exact \\
strength of the warped coupling over the flat one & 1 & $1.8\times10^{-16}$ & $1\times10^{-12}$ \\
Goldberger--Wise mass coefficient: $\frac43$ and the note's $\frac23$ & 2 & 0 & $1\times10^{-15}$ \\
Higgs--radion: $c=-6$, $K_{hr}/(\xi\gamma)$, $(K_{rr}-1)/(\xi\gamma^2)$ & 3 & 0 & exact \\
Higgs--radion mixing in 7+1: $Z-1$, $\theta$, rate changes; warped $\gamma$ & 5 & $4.4\times10^{-16}$ & $1\times10^{-12}$ \\*
\end{longtable}}
\Needspace{8\baselineskip}\subsection*{6.3 The earlier sections are unchanged}\addcontentsline{toc}{subsection}{6.3 The earlier sections are unchanged}
tools/verify\_7p1\_nb.py compares the evaluated notebook with the backups of section 3 and with the Python programs; its lines on sections 25 and 26:

\begin{Verbatim}[fontsize=\small,frame=single,framesep=2mm]
section 25: check outputs (all True) : 27 (True)
stage 6: without section 25 (and with the old title of the summary section) the 786 input and
    ...text cells equal those of the backup taken before section 25 was appended : True
stage 6: the 481 outputs of sections 0-24 and of the summary equal those of the backup (summary
    ...counts and tables, the private context name and measured times excepted) : True
section 25 numbers: 17 groups of printed numbers agree with tools/ref_torus_brane.py (integers
    ...and exact values exactly, floating-point numbers to 1e-15 ... 1e-6 per group) : True
section 26: check outputs (all True) : 33 (True)
stage 7: without section 26 (and with the old title of the summary section) the 815 input and
    ...text cells equal those of the backup taken before section 26 was appended : True
stage 7: the 515 outputs of sections 0-25 and of the summary equal those of the backup (summary
    ...counts and tables, the private context name and measured times excepted) : True
section 26 numbers: 15 groups of printed numbers agree with tools/ref_radion.py (integers and
    ...exact values exactly, floating-point numbers to 1e-15 ... 1e-9 per group) : True
VERDICT: ALL CROSS-CHECKS PASSED
\end{Verbatim}
\Needspace{8\baselineskip}\subsection*{6.4 Files}\addcontentsline{toc}{subsection}{6.4 Files}
\textbf{SHA-256 of the files} (at the time this document was generated):

\begin{Verbatim}[fontsize=\small,frame=single,framesep=2mm]
7+1.nb                             41d252f9852f210c228f45e1ac956d2d7ef26faeb0062a5880280e1aae0ef
    ...fc7
tools/make_7p1_nb.py               996d26c9552cd44902560af16a1683fa71457261f18fefaeddae87bd99a2b
    ...144
tools/nb_part_gravity.py           c4caf2ad5c932e8f603341edade5f4b98fa784feb05906c051eab7c884cd6
    ...8f8
tools/nb_part_ks_gravity.py        a656c0ad040b4576b75b17e029545109e531c8e656ad99bfdb5cba946e967
    ...772
tools/nb_part_sm.py                43e92beed57453a4b57931de6d603ac6c218efba6e4005133b69b8054a975
    ...c5b
tools/nb_part_coupled.py           e630521d64132e71afc861114f699eb06454557cf14a57ad8ec1eef96247a
    ...ae0
tools/nb_part_dark.py              10ad8fb8ae25fa07bbee758c653995acd4e2c80733059a8b3648bd1593e46
    ...db0
tools/nb_part_torus_brane.py       fb904ff8d85b904c55b3b2a8f60fb0cb047f1ccf1b98db3364092f2bae0b2
    ...162
tools/nb_part_radion.py            53c5604707a3e22228c840ee0a2d2125616f9410f75a5535041dfaf18f497
    ...c00
tools/lovelock71.py                d93ecbd0f82adb972218d59dfe32999a4af944945ca9758607f3a633e1af9
    ...5a1
tools/ref_gravity.py               f07cf8adaff585af89c87d19532f1aaa7db606ad88ae90483b2d1f41341ed
    ...f36
tools/ref_ks_gravity.py            b6488ef24b1c16aacd8ff6cf6a5a1ba3f316e082449ad97480022d1a26383
    ...27f
tools/ref_sm.py                    f3f80806893f42bbf166e1f1de896ffd611dbfe847685da8e3ee367e9cde6
    ...aba
tools/ref_coupled.py               695dbd9375753653472d9c37502f0067c4a64fe01a64ba8c05e8b2c5da85a
    ...a58
tools/ref_dark.py                  1600e6b57ab0383b4f2b664e2796d3864c9afa77a9ed59f007f66e1139453
    ...ff5
tools/ref_torus_brane.py           1f623656ab0f5c8ee43cb310f7bfd580cbba5a7385da2963adaae4d2f8bda
    ...af7
tools/ref_radion.py                e9d47176b9f2d2232a84b3b82188bd3b0a19cde3dbe1ba013c1fe5abff634
    ...1ce
tools/ref71.py                     1a7a06eb0add3f459611776dc9c474a5bb959d947a9db949db4612b1ab943
    ...4e0
tools/ref44.py                     ab100c3ebcc055fbe59c132a02f80ecb17630e934e99d59f88e919ec8064f
    ...af5
\end{Verbatim}
\end{document}
